% ============================================================
%  第十一章后半：对读者提问的回应（§11.4 至章末）
%  本文件由 ch11a.tex 之后 \input 拼接，直接以节标题开始
% ============================================================

\providecommand{\qhead}[1]{\par\addvspace{0.7em}\noindent{\heiti #1}\par\nopagebreak\addvspace{0.15em}}
\providecommand{\said}[1]{\par\noindent\textbf{#1}\quad}

\section{政策评估与 do 算子}

\subsection{识别条件计划（\S4.2，第 113 页）}

\qhead{问题（致作者）}

本书 \S4.2（第 113 页）给出了条件行动之效应的识别条件与估计公式，即行动 $\dox{X=g(z)}$ 的效应，其中 $Z=z$ 是在行动之前取得的一项测量。该式能否推广到多个行动的情形，即“条件计划”（conditional plan）？

我们看到的困难在于：该公式是在“$X$ 不改变 $Z$ 的取值”这一假设下导出的。然而在包含多个行动的计划中，$X$ 中的某些行动可能改变那些指导后续行动的观测 $Z$。我们目前没有记号来区分干预后的观测与干预前的观测。缺少这种记号，条件计划如何被形式地表达、并提交给 do 演算分析，就不清楚。

\qhead{作者答复（与 Ilya Shpitser 合作）}

干预后观测与干预前观测之间的记号区分，由第七章借助反事实语言引入。不过，条件计划的情形无须借助更丰富的记号即可处理。理由在于：决定某一行动之选择的那些观测，不会被该行动本身改变；而已经被先前行动改变的那些观测，则由 $P(y\mid\dox{x},z)$ 这一记号充分刻画。

为了看清这一点，我们先引入反事实记号，然后证明它可以从所得表达式中消去。我们用粗体字母表示集合，用普通字母表示集合中的单个元素；大写字母表示随机变量，小写字母表示这些变量可能取得的取值。我们写 $Y_x$，意思是“若把变量 $X$ 设定为取值 $x$，$Y$ 所取得的取值”。类似地，$Y_{X_g}$ 意指“若把变量 $X$ 设定为它们在随机策略 $g$ 之下本会取得的那些取值，$Y$ 所取得的取值”。注意 $Y_x$ 与 $Y_{X_g}$ 都是随机变量，与原变量 $Y$ 无异。

设有 $K$ 个按某种时间顺序发生的行动变量 $X$。我们用上标标明某个变量被施以行动的时点，于是当 $i<j$ 时，变量 $X^i$ 先于 $X^j$ 发生。对给定的 $X^i$，记 $X^{<i}$ 为 $X^i$ 之前的行动变量集合。

我们所关心的是：在一项策略之下，一组结局变量 $Y$ 的概率分布。该策略把每个 $X^i\in X$ 的取值设定为函数 $g_i$（事先已知）的输出，而这些函数依据某组在先变量 $Z^i$ 以及此前对 $X^{<i}$ 的干预取值。与此同时，变量 $Z^i$ 本身又受先前干预的影响。为恰当地定义这一递归，我们采用归纳定义。基础情形是 $X^1_g=g_1(Z^1)$；归纳情形是
\[
X^i_g=g_i\bigl(Z^i_{X_g^{<i}},\,X^{<i}_g\bigr).
\]
这里下标 $g$ 表示我们所使用的策略，即 $g=\{g_i\mid i=1,2,\dots,K\}$。于是可以写出所关心的量：
\[
P(Y_{X_g}=y)=P\bigl(Y=y\mid \dox{X^1=X^1_g},\dox{X^2=X^2_g},\dots,\dox{X^K=X^K_g}\bigr).
\]
令 $Z_g=\bigcup_i Z^i_{X_g^{<i}}$。关键观察是：一旦观测到 $Z_g$ 取定某些值，$X_g$ 也就塌缩为唯一的值，因为 $X_g$ 是 $Z_g$ 的函数。记
\[
x_z=\{x^1_z,\dots,x^K_z\}
\]
为 $X_g$ 在“$Z_g$ 被观测为 $z=\{z^1,\dots,z^K\}$”这一情形下所取得的取值。这里要指出：若已知 $z$，我们便可事先算出 $x_z$，因为函数 $g_i$ 是事先固定且为我们所知的。然而我们并不知道 $Z_g$ 会取得哪些值，于是用分情形分析遍历一切可能的取值组合，得到
\begin{align*}
P(Y_{X_g}=y)&=\sum_{z^1,\dots,z^K}P\bigl(Y=y\mid\dox{X=x_z},Z^1=z^1,\dots,Z^K=z^K\bigr)\\
&\qquad\times P\bigl(Z^1=z^1,\dots,Z^K=z^K\mid\dox{X=x_z}\bigr).
\end{align*}
这里我们注意到，$Z^i$ 不可能依赖于其后的干预，于是得到
\begin{align*}
P(Y_{X_g}=y)&=\sum_{z}P\bigl(Y=y\mid\dox{X=x_z},Z^1_{x_z}=z^1,\dots,Z^K_{x_z}=z^K\bigr)\\
&\qquad\times P\bigl(Z^1=z^1,\dots,Z^K=z^K\mid\dox{X=x_z}\bigr).
\end{align*}
再注意到第一、第二两个因子中的下标已是多余的，因为 $\dox{x_z}$ 已为该表达式中一切变量隐含了这类下标。于是目标量可改写为
\begin{align*}
P(Y_{X_g}=y)&=\sum_{z}P\bigl(Y=y\mid\dox{X=x_z},Z^1=z^1,\dots,Z^K=z^K\bigr)\\
&\qquad\times P\bigl(Z^1=z^1,\dots,Z^K=z^K\mid\dox{X=x_z}\bigr),
\end{align*}
或者更简洁地写作
\[
P(Y_{X_g}=y)=\sum_{z}P\bigl(y\mid\dox{x_z},z\bigr)P\bigl(z\mid\dox{x_z}\bigr).
\]
可见，只要 $Y$、$X$、$Z$ 是互不相交的集合，该式便可由 $P(y\mid\dox{x},z)$ 与 $P(z\mid\dox{x})$ 算出。

总结一下：条件计划虽由嵌套的反事实表达式自然表示，其识别问题却仍可归约为形如 $P(y\mid\dox{x},z)$ 的条件干预分布的识别问题（$z$ 可以为空）。在给定图 $G$ 中从联合分布识别这类量的完整条件，由 Shpitser and Pearl (2006a,b) 给出。

\subsection{间接效应的含义}

\qhead{问题（致作者）}

我在讲授一门潜变量建模课程（面向生物统计及其他公共卫生专业的学生），昨天正介绍通径分析的概念，其中包括直接效应与间接效应。

我演示了如何以直接路径之积来计算间接效应。随后有学生问该如何解释间接效应，我给出了我一贯给出的回答：间接效应 $ab$（在图 11.13 的简单模型中）是 $X$ 的变化经由其与 $Z$ 的关系而对 $Y$ 产生的效应。

\begin{figure}[htbp]
\centering
\begin{tikzpicture}[causal,scale=1.0,every node/.style={transform shape}]
  \node[circle,fill,inner sep=1.4pt] (X) at (0,1.0) {};
  \node[font=\small,left=0.6mm] at (0,1.0) {$X$};
  \node[circle,fill,inner sep=1.4pt] (Y) at (2.6,1.0) {};
  \node[font=\small,right=0.6mm] at (2.6,1.0) {$Y$};
  \node[circle,fill,inner sep=1.4pt] (Z) at (1.3,-0.9) {};
  \node[font=\small,below=0.6mm] at (1.3,-0.9) {$Z$};
  \draw[ed] (X)--(Y) node[midway,above=0.6mm,font=\small] {$c$};
  \draw[ed] (X)--(Z) node[midway,left=0.8mm,font=\small] {$a$};
  \draw[ed] (Z)--(Y) node[midway,right=0.8mm,font=\small] {$b$};
\end{tikzpicture}
\caption{演示 $X$ 经由 $Z$ 对 $Y$ 的间接效应。}
\label{fig:11.13}
\end{figure}

略一沉吟之后，学生问了下面这番话：

\said{学生：}“$b$ 这条路径的解释是：$b$ 是在固定 $X$ 的条件下，$Z$ 每增加一个单位时我们在 $Y$ 上看到的增量，对吗？”

\said{我：}“对。”

\said{学生：}“那么，当我们解释间接效应时，被固定不变的又是什么？”

\said{我：}“不太明白你的意思。”

\said{学生：}“您说间接效应 $ab$ 的解释是：$ab$ 是在 $X$ 增加一个单位、经由其对 $Z$ 的因果效应时，我们在 $Y$ 上看到的增量。可是既然 $b$（从 $Z$ 到 $Y$ 的直接效应）要求固定 $X$，它怎么能被用在一个同时要求 $X$ 变动一个单位的计算里呢？”

\said{我：}“嗯，问得很好。我恐怕没有好答案给你。在从 $X$ 到 $Y$ 的直接路径为零的情形下，我想没有问题，因为那时 $Z$ 与 $Y$ 的关系与 $X$ 无关。但你说得对，此处若 $c$ 非零，我们就必须把 $b$ 解释为固定 $X$ 时 $Z$ 对 $Y$ 的效应。我明白这听起来与 $ab$ 间接效应的解释相冲突——后者考察的是 $X$ 的变化会引起什么。这样吧，我回头再答复你。我早跟你说过，这里的计算并不难，难的是真正理解你的模型意味着什么。”

\qhead{作者答复}

请代我称赞这位学生好问的精神。答案其实可以讲得相当简单（见第二版所补写的 \S4.5.5）：

\emph{$X$ 对 $Y$ 的间接效应，是在固定 $X$ 不变、同时把 $Z$ 提高到“$X$ 增加一个单位时 $Z$ 本会取得的那个值”的条件下，我们在 $Y$ 上看到的增量。}

这一反事实定义引出中介公式 \textup{(4.18)}，它把通径分析推广到非线性模型，并使我们能用普通回归来估计类别变量的间接效应。

\qhead{作者补记}

这个问题代表了结构方程模型（SEM）的标准教育亟待改革的若干方面之一。虽说有些 SEM 教材会顺带提到结构参数可解释为效应系数，但作者、教师与学生都并不认真对待这一解释。写到这里已是 2008 年，我发现 SEM 教育的主体仍然聚焦于统计估计与模型拟合的技术，而一旦模型拟合、估计完毕，几乎找不到关于“模型意味着什么”的严肃讨论（SEM 求生指南见 \S11.5.3）。\footnote{在我所考察过的 2000 年以后出版的 SEM 教材中，没有一本的索引里出现“causal”一词。}

这一教育传统的弱点，在好问的学生提出略微偏离标准 LISREL 套路的问题时便暴露出来——这些问题的答案，取决于对结构系数与结构方程的因果解释。例如：

\begin{enumerate}[label=\arabic*.,leftmargin=1.8em]
\item 我们为什么要以那样的方式定义总效应（即定义为若干直接效应之积的和）？这是一个任意的定义，还是由路径系数的因果解释所迫？
\item 我们为什么要把间接效应定义为总效应与直接效应之差？
\item 在非线性系统中，或在含二分变量的系统中，我们如何定义直接效应？
\item 在含反馈环（即互为因果）的系统中，我们应如何以有意义的方式定义效应，从而避开 SEM 教材所引用的那些错误定义带来的陷阱？（见第 164 页）
\item 如果我们试图估计某一行动的效应，并在实施该行动之前先作一些测量，那么我们对直接效应与总效应的评估是否仍然不变？
\end{enumerate}

读者会乐于知道，这些问题都有形式化的答案，见 \S4.5.4、\S4.5.5、\S11.5.2、\S11.5.3 与 \S11.7.1。

就个人而言，我对直接效应与间接效应的兴趣，是由 Jacques Hagenaars 的一封来信激起的。他在信中写道（2000 年 9 月 15 日）：“间接效应在实质理论中确实占据重要位置。许多社会科学理论在输入变量（背景特征）与输出变量（行为）上‘意见一致’，分歧恰恰在于中间机制。举个简单的例子：我们知道教育对政治偏好的影响，一部分经由‘经济地位’中介（受教育程度高的人得到更好的工作、挣更多的钱），一部分经由‘文化机制’中介（与教育内容以及学校里伴随的社会化过程有关）。我们需要认识并区分这两种不同过程的性质与后果，也就是说，我们想知道间接效应的符号与大小。在结构方程模型的参数化线性版本中，存在一套‘路径系数演算’，我们可以用它把总效应写成直接效应与若干间接效应的表达式。但在一般的非参数情形下这做不到，例如在对数线性参数化版本中也做不到。至于逻辑模型系统，早已有人指出并不存在可与之相比的‘路径系数演算’。然而，鉴于其压倒性的理论重要性，间接效应问题不能就这样被忽略。”

受这些评论激励，又有了嵌套反事实的记号，我便着手把第 147 页给出的“雇用歧视”的法律定义形式化，并由此得到 \S4.5.4 与 \S4.5.5 的结果——其中一部分已由 Robins and Greenland (1992) 预见——以及式 \textup{(4.12)}、\textup{(4.17)} 与 \textup{(4.18)} 的中介公式。受这些结果的启发，我不得不、也很乐意收回自己在《因果性》第一版第 165 页所下的断言：“间接效应缺乏内在的操作含义”，理由是它们无法用 $\dox{x}$ 算子分离出来。间接效应确实无法用 $\dox{x}$ 算子分离，但它们确实具有内在的操作含义。直接效应与间接效应的政策含义，在 Pearl (2001, 2010b) 与 Petersen 等 (2006) 中得到了进一步例示。

\subsection{do(x) 能表示实际的实验吗？}

\qhead{问题（致作者）}

亚利桑那大学的 L.B.S. 质疑 $\dox{x}$ 算子能否表示现实的行动或实验：“即便是一个在其他方面执行得完美的随机化实验，也可能得出完全误导的结论。一个很好的例子是一项研究：以注射维生素 E 作为疗法，用于治疗有晶状体后纤维增生症风险的暖箱婴儿。该随机化实验显示了注射有效，但人们不久便发现，真正起作用的疗法是每天数次打开加压的、氧饱和的暖箱以便注射，从而降低了婴儿血液中的气压与氧水平（Leonard, \textit{Major Medical Mistakes}）。在那种情况下，任何统计分析都会是误导的。”

佐治亚理工学院的 S.M. 补充道：“你那个误导性因果效应的例子，正显示了让我对 $\dox{x}$ 概念感到不安的那类事情。你要么 $\dox{x}$、要么不 $\dox{x}$，但真正的原因可能是某个与 $\dox{x}$ 共变的别的东西，而不是 $\dox{x}$ 本身。”

\qhead{作者答复}

数学处理的是理想情形，而实验者的职责，是确保实验条件尽可能接近数学理想。$\dox{x}$ 算子代表的是在一个理想实验中实施 $X=x$：其中被直接操纵的只有 $X$、且仅仅是 $X$，而不是模型中的任何其他变量。

在你举的维生素 E 注射一例中，还有另一个变量与 $X$ 一起被操纵了，即暖箱盖 $Z$；这就把实验变成了 $\dox{x,z}$ 条件，而不是 $\dox{x}$。因此，该实验远非所设想的那个实验，甚至远不符合标准的实验规程——后者要求使用安慰剂。若使用了安慰剂（以逼近 $\dox{x}$ 算子的要求），结果就不会有偏。

除非指明模型中有哪些变量被该行动直接影响、以及如何影响，任何模型都无法预测一个行动的效应。$\dox{x}$ 算子是一种数学装置，它帮助我们形式地指明：在任一给定实验中，什么被保持恒定、什么可以自由变化。随后，do 演算帮助我们推演这类设定的逻辑后果——前提是这些设定被忠实执行，前提是我们拥有关于环境的有效因果模型，并且我们拥有在其他条件得到良好规定的实验中收集到的数据。参见 \S11.4.6。

\subsection{do(x) 算子具有普适性吗？}

\qhead{问题（致作者，来自 Bill Shipley）}

在大多数实验中，外部操纵在于给 $X$ 增加（或减去）某个量，而不移除 $X$ 既有的原因。例如，给一块田地增施 5 kg/h 的肥料，给受试者增加 5 mg/l 的胰岛素，等等。这里，被操纵变量既有的原因仍在起作用，只是新增了一个变量（$M$）。

我认为，把 $\dox{x}$ 算子当作外部操纵的一般算子，问题在于它要求两件事：(1) 移除 $x$ 的一切既有原因；(2) 把 $x$ 设定为某个值。这对应于某些类型的外部操纵，但并不对应于全部（甚至大部分）外部操纵。我愿引入一个 $\mathrm{add}(x+n)$ 算子，意思是“在既有的因果过程之外，给 $x$ 增加量 $n$”。

在图形上，这相当于在既有的因果图上增补一条新边，即 $M\to X$（其中 $M=n$）。在代数上，这相当于把一项新的项 $n$ 作为 $X$ 的原因加入（Shipley 2000b）。

\qhead{作者答复}

在许多情形下，你所说的“加性干预”确实是我们能对变量 $X$ 施加干预的唯一方式；在另一些情形下，它可能正是我们希望评估的实际政策。事实上，干预的一般概念（第 113 页）允许我们把 $X$ 的方程替换为任何其他适合具体情形的方程，而不必是常数 $X=x$。你所提议的做法，相当于把 $X$ 的旧方程 $x=f(\pa_X)$ 替换为新方程 $x=f(\pa_X)+n$。这一替换可用“工具变量”来表示，因为它等价于写出 $x=f(\pa_X)+I$（$I$ 为工具），并让 $I$ 从 $0$ 变到 $n$。

有三点值得注意：

\begin{enumerate}[label=\arabic*.,leftmargin=1.8em]
\item 加性操纵可以在 $\dox{\cdot}$ 框架内表示——我们只是把 $\dox{\cdot}$ 算子施加于工具 $I$，而不施加于 $X$ 本身。这是一种不同种类的操纵，需要与 $\dox{x}$ 相区别，因为如下文所示，它对 $Y$ 的效应可能不同。
\item 在许多情形下，科学家并不满足于估计工具对 $Y$ 的效应，而是努力估计 $X$ 本身的效应，因为后者往往更有意义，也更容易迁移到其他情境。（关于“意向治疗”（intention to treat）效应的讨论见第 261 页，例子见第 363 页。）
\item 考虑 LISREL 失效的那个非递归例子（第 164 页）。若把“总效应”解释为 $Y$ 对工具 $I$ 一个单位变化的响应，则得到 LISREL 的公式：$I$ 对 $Y$ 的效应为 $b/(1-ab)$。然而，若我们坚持“$X$ 每变化一个单位”的概念，则回到 do 公式：$X$ 对 $Y$ 的效应是 $b$，而不是 $b/(1-ab)$——尽管操纵是借助工具完成的。换言之，我们把 $I$ 从 $0$ 变到 $1$，观察 $X$ 与 $Y$ 的变化；若以 $X$ 的变化去除 $Y$ 的变化，得到的是 $b$，而不是 $b/(1-ab)$。
\end{enumerate}

总结一下：是的，加性操纵有时正是我们需要建模的东西，而且它可以在 $\dox{x}$ 框架内借助工具变量完成。不过，我们仍须区分工具的效应与 $X$ 的效应：前者不稳定（第 261 页），后者稳定。Shpitser and Pearl (2009) 给出了识别“给 $X$ 加 $n$”算子之效应的充要图形条件。

\qhead{Bill Shipley 进一步追问}

多谢澄清。在我看来，对因果系统的操纵进行建模与表示，最简洁、最直截了当的方式就是：(1) 修改未操纵系统的因果图，以表示所提议的操纵；(2) 把这张新图翻译为结构方程；(3) 从所得方程推出预测（包括条件预测）。我在自己的书里就是这样处理这一概念的。既然如此，何必还要操心 $\dox{x}$？尤其是，人们可以用这种方式对相当精巧的操纵建模。例如，完全可以问：若给因果图中的某个变量 $x$ 增加量 $z$，而 $z$ 又依赖于图中的另一个变量，会发生什么？

\qhead{作者答复}

你所提议的方法，是把当前的方程 $x=f(\pa_X)$ 替换为 $x=g\bigl(f(\pa_X),I,z\bigr)$；它要求我们知道 $f$ 与 $g$ 的函数形式（如在线性系统中那样），或者要求 $X$ 的父节点可观测（如第 74 页的过程控制一例那样）。然而，在第三章与第四章的非参数、部分可观测设定中，这两条都不成立；这就可能使得：从干预之前收集的数据出发，无法预测所提议干预的效应——我们把这一问题称为识别。由于变量 $I$ 的干预前统计量不可得，而 $f$ 又未知，因此存在这样的半马尔可夫情形：$P(y\mid\dox{x})$ 可识别，而 $P\bigl(y\mid\dox{x=g(f(\pa_X),I,z)}\bigr)$ 不可识别；每个情形都必须按其自身的是非曲直来分析。因此，重要的是为这一大片潜在的干预空间设定某些标准，并把注意力集中在那些能够照亮其他干预的干预上。

科学靠标准而繁荣，因为标准至少有双重作用：交流与理论聚焦。举例来说，数学家认定导数算子“$dy/dx$”是交流变化信息的好标准，于是我们在微积分课里教它，尽管别的算子也能起同样作用，例如 $x\,dy/dx$ 或 $(dy/dx)/y$ 等等。因果分析亦然：

\begin{enumerate}[label=\arabic*.,leftmargin=1.8em]
\item \textbf{交流}：倘若我们从流行病学中取消“处理效应”一词，代之以对效应如何被测量所作的详尽描述，我们实际上会扼杀流行病学家之间的一切交流。于是确立了一个标准：在受控的随机化实验中所测量的东西，称为“处理效应”；其余则视为该主题的变奏。“do 算子”忠实地代表了这一标准。SEM 亦然。Sewall Wright 谈论“效应系数”，并使之成为通径分析中“直接效应”的标准（那是在它被回归学的行话搅乱之前）；借助它，可以构造更精细的效应。同样，“do 算子”是定义这一标准的基础。
\item \textbf{理论聚焦}：操纵的许多变体都可以化归为“do”，或化归为若干次“do”的施用。于是，为“do”所建立的理论结果便可适用于这些变体。例如：你的“加 $n$”操纵可表达为对工具施加“do”。又如：涉及“do”的表达式的识别问题，可适用于更精细效应的识别问题。例如第 113 页表明，若表达式 $P(y\mid\dox{x},z)$ 可识别，则条件行动 $P\bigl(y\mid\dox{x=g(z)\ \text{若}\ Z=z}\bigr)$ 的效应也可识别。书中许多其他理论结果亦然：它们是为“do”算子发展的，彼此借鉴，并且适用于许多变体。
\end{enumerate}

最后，do 算子背后的“手术”操作，为解释反事实语句提供了恰当的形式体系（见第 204 页），而反事实在科学话语中比比皆是（见第 217—219 页）。我至今尚未见到任何竞争者能在通用性、普遍性、形式力量以及（并非最不重要的）概念吸引力上与之相比。

\subsection{没有操纵的因果！！！}

\qhead{问题（致作者）}

在分析直接效应时，\S4.5 举了学校录取中的性别歧视一例，并在若干公式中把性别置于“帽”记号之下，或等价地作为 do 算子的一个变元。性别是一个无法被操纵的变量，怎么能放在 do 算子之后呢？

\qhead{作者答复}

自从 Holland 提出“没有操纵就没有因果”（Holland 1986）这一口号以来，许多好想法在因果分析中被扼杀或弃置了。禁止谈论性别如何造成男女之间在生物学、社会与心理上的诸多差异，等于压制我们关于性别差异的知识的 90\%。

我们当然拥有没有操纵的因果。月亮造成潮汐，种族造成歧视，性别造成某些激素而非另一些激素的分泌。自然界是一个由机制组成的社会，这些机制不倦地感知某些变量的取值并决定另一些变量的取值；它并不等待人类操纵者来启动这些机制。

诚然，操纵是科学家检验机制运作的一种方式（尽管是粗糙的一种），但它绝不应妨碍对推动所研究现象之机制的因果思考、形式定义与数学分析。也许正因如此，科学家才发明了反事实：它使他们能够陈述并设想前件条件的实现，而不必指明这些条件赖以建立的具体物理手段。

定义 4.5.1 中“帽”记号的目的，不是要激发关于改变申请人性别之可能途径的思考，而是要提醒我们：任何涉及“效应”的定义都应聚焦于因果联系，并把虚假关联从所定义的数量中过滤掉。诚然，在性别一例中，我们可以放心地把 $P(y\mid \text{女})$ 替换为 $P(y\mid\dox{\text{女}})$，因为决定性的机制可以稳妥地假定为独立于影响 $Y$ 的背景因素（从而保证没有混杂）。但作为一个一般定义，甚至作为教学例子的一部分，涉及直接效应与性别歧视的数学表达式都应保留帽记号。至少，把“女”置于帽记号之下，有助于传播那句早该出现的反口号：“没有操纵的因果？当然有！”

\subsection{与 Cartwright 一起猎取原因}

Nancy Cartwright 在其《猎取原因并加以利用》（Hunting Causes and Using Them，剑桥大学出版社，2007）一书中，对 $\dox{x}$ 算子以及它所依据的“手术”（surgery）语义表达了若干反对意见（第 72 页、第 201 页）。她在这样做时，揭示了若干亟需系统澄清的领域；我将逐一回应。

Cartwright 对手术的描述如下：

\begin{quote}\kaishu
Pearl 为反事实给出了一套精确而细致的语义。但这套语义是关于什么的语义？Pearl 所发展的这套特定语义，并不适合反事实在自然语言中的大量用法，尤其不适合我一直在讨论的、用于规划与评估的那类用法。原因在于，他想象反事实前件被实现的那种特殊方式：通过一次精确的切口，恰好改变反事实前件而不改变任何别的东西（除了由这一差别因果地引出的东西）。但当我们考虑实施一项政策时，这根本不是我们需要问的问题。对于政策与评估，我们通常想知道的是：假如政策真的被付诸实施，会发生什么。而无论我们对它可能如何被付诸实施知道些什么，我们通常能确定的唯一一件事是：它不会以 Pearl 所设想的那种精确切口的方式实现。

以 Pearl 的复合公理为例，他证明该公理在所有因果模型中都成立——给定他对因果模型的刻画以及他对反事实的语义。该公理称：“若我们迫使一个变量（$W$）取它在没有我们干预的情况下本会取得的取值 $w$，那么该干预对系统中其他变量没有影响”（第 229 页）。如果我们设想的是以尽可能小的方式实现反事实前件的干预，这条公理是合理的。但在大量现实情形中，它显然被违反。我们常常根本不知道前件是否真的会实现，即便我们承认支配性原理是确定性的，情况亦然。我们实施一项政策以确保前件实现——而该政策可能在系统中引起其他变量的一大堆变化，有些是预料之中的，有些则不是。（Cartwright 2007, pp.~246--7）
\end{quote}

于是，Cartwright 的反对意见可概括为三条主张，下面逐一回应。

\begin{enumerate}[label=\arabic*.,leftmargin=1.8em]
\item 在大多数研究中，我们需要预测非原子干预的效应。
\item 对于政策评估，“我们通常想知道的是：假如政策真的被付诸实施，会发生什么”，但不幸的是，“该政策可能在系统中引起其他变量的一大堆变化，有些是预料之中的，有些则不是”。
\item 由于实际政策是非原子的，即便我们能够设想出受政策影响的那些变量，也无法用 $\dox{x}$ 演算的原子语义来评估它们。
\end{enumerate}

先从主张 (2) 说起——它最容易驳倒。这一反对意见与 \S11.4.3 所讨论的那条完全相同，其答案是：“除非正确指明模型中有哪些变量受该行动影响、以及如何影响，任何模型都无法预测一个行动的效应。”换言之，在 Cartwright 主张 (2) 所描述的那种无知状态下，政策评估研究只能以一个平淡的答案收场：信息不足，因而什么都可能发生。这就像在黑暗中按下一个不知用途的按钮，或试图用两个方程解三个未知数。更何况，do 演算本身就可用来检验：在任一给定情境中，这种无知状态是否足以使那个平淡的答案成为正当的。因此，若以为严肃的政策评估研究是在这种无知状态下进行的，那就错了；我所见过的所有政策分析，都始于假定我们知道哪些变量受政策影响，并把这一知识形式地表达出来。

主张 (1) 或许适用于某些情形，但肯定不会适用于大多数情形；在许多研究中，我们的目标不是预测我们即将实施的那种粗糙的非原子干预的效应，而是评估一项理想的原子政策——它在现有工具条件下无法实施，却代表了一种对我们的领域理解至为关键的理论关系。

一个例子会有帮助。以我们今天可用的任何法律或教育手段，都无法使吸烟停止；但香烟广告可以。这并不妨碍研究者致力于估计“吸烟对癌症的效应”，并且是从那些改变工具——香烟广告——而非改变吸烟的实验中去估计。

他们之所以对原子干预 $P(\text{癌症}\mid\dox{\text{吸烟}})$ 感兴趣，而不是（或不止于）对 $P(\text{癌症}\mid\dox{\text{广告}})$ 感兴趣，是因为前者代表该人群的一项稳定的生物学特征，不受那些影响广告敏感性的社会因素的污染。借助这一稳定特征，人们可以评估各种各样的实际政策，而每项政策都使用不同的减烟工具。

最后，主张 (3) 在本书几乎每一章中都被确凿地驳倒了。还有什么比一项涉及一序列行动的政策更“非原子”呢？其中每个行动都由一组观测 $Z$ 决定，而这些观测又受先前行动的影响（见 \S4.4 与 \S11.4.1）。然而，实施这样一项复杂政策的效应，却可以用 do 演算的“手术”语义来预测，其方式与用原子物理学预测复杂分子的性质大体相同。

我曾向 Nancy Cartwright 提出挑战（Pearl 2003a），现在愿再挑战一次：请举出一个政策实例，它要么无法用 $\dox{x}$ 算子加以规定与分析，要么无法同样用 $\dox{x}$ 算子演算证明其“不可预测”（例如在黑暗中按下不知用途的按钮）。

具有讽刺意味的是，回避基于理想原子干预的数学，反倒可能使科学家在处理现实的非原子干预时笨拙无能。

科学与数学中充满了辅助性的抽象量：它们不被直接测量或检验，却服务于对可测量的分析。纯净的化学元素在自然界中并不存在，然而对于理解合金与化合物不可或缺。负数（更不用说虚数）并不孤立存在，然而对于理解对正数的运算至关重要。

本书所处理（并解决）的广泛问题表明：关于干预与实验的种种问题——理想的与非理想的、实践的和认识论的——无一例外都可以用原子干预作为原初概念来精确表述并系统处理。

\subsection{非模块性的幻象}

在批判 do 算子时，Cartwright 还提出了另一条论据——模块性的失效，据称这一失效困扰着大多数机械系统与社会系统。

用她的话说：

\begin{quote}\kaishu
“Pearl 最近在 LSE 谈到这一点时，用一个电路的布尔输入—输出图来说明这一要求。在那张图里，不仅每个变量的整个输入可以独立于其他变量的输入而改变，而且该输入的每个布尔分量也可以独立改变。但我们研究的大多数装置并非如此。它们更像一台烤面包机或一个化油器。”
\end{quote}

至此，Cartwright 给出了一个汽车化油器的四方程模型，并总结道：

\begin{quote}\kaishu
腔内的气体是泵入的汽油与流出乳化管的汽油之和。二者各贡献多少，由其他因素固定：对泵入的汽油而言，是空气流量和一个参数 $a$，后者部分地由腔体的几何形状决定；对流出乳化管的汽油而言，则是一个同样取决于腔体几何形状的参数 $a'$。要点在于：在 Pearl 的电路板中，每一条不同的因果联系都有各自不同的物理机制来支撑。但那对空间与材料是惊人的浪费，而这对化油器来说是要紧的。工程师设计化油器的核心技巧之一，就是确保同一种物理设计——例如腔体的设计——能够同时支撑或保证我们需要的若干不同因果联系。

只要回头看我的图解方程，就能看到大量定律都依赖于同一批物理特征——化油器的几何形状。所以这些定律中没有一条可以被单独改变。要改变其中任何一条，都需要重新设计化油器，而这会连带改变其他定律。按设计，不同的因果定律被捆绑在一起，无法单独改变。于是模块性失效了。（Cartwright 2007, pp.~15--16）
\end{quote}

因此，在 Cartwright 看来，一组共享参数的方程本质上就是非模块的；改变其中一个方程意味着至少改动它的一个参数，而若该参数出现在别的方程中，它也必须随之改变，这就违反了模块性。

Heckman (2005, p.~44) 提出了类似的主张：“对一个方程施加约束，就是对整个内生变量集合施加限制。”“关掉一个方程，也可能影响系统中其他方程的参数，从而违反参数稳定性的要求。”

这类恐惧与警告都是幻象。手术以及围绕它建立起来的整套语义与演算，并不假设在物理世界中我们拥有这样的技术：能够精确地改动每个结构方程背后的机制，同时让其余一切保持不变。符号层面的模块性并不假设物理层面的模块性。手术是一种符号操作，它对实验者可用的物理手段不作任何断言，对相关机制之间可能存在的不可见联系也不作任何断言。

在符号层面，人们当然可以改变一个方程而不触动其他方程，并进而定义依赖于这类“原子”改变的量。这样定义的量是否对应于能够在物理上实现的变化，完全是另一个问题，只有在我们对可用的干预有了形式描述之后才能回答。更重要的是，关掉一个方程并不必然意味着摆弄它的参数；它意味着否决该方程，即让方程保持原样，却把结局变量从其影响下解放出来。

一个简单的例子可以说明这一点。

设有两个处于自由落体状态下的物体，其加速度分别为 $a_1$ 与 $a_2$，由方程
\begin{equation}
a_1=g \tag{11.12}
\end{equation}
\begin{equation}
a_2=g \tag{11.13}
\end{equation}
给出，其中 $g$ 是地球引力。两个方程共享参数 $g$，在 Cartwright 的意义上显得是非模块的；要把引力作用于一个物体而不相应地改变对另一物体的引力，确实没有物理办法。然而，这并不意味着我们不能在不触动物体 2 的情况下对物体 1 施加干预。设想我们抓住物体 1 并使之停住。在数学上，这一干预相当于把式 \textup{(11.12)} 替换为
\begin{equation}
a_1=0 \tag{11.14}
\end{equation}
而让式 \textup{(11.13)} 保持原样。在式 \textup{(11.12)} 中把 $g$ 置零，是一次符号手术：它不改变物理世界中的 $g$，而是通过让物体 1 受到一个来自我们抓握之手的新的力 $f$ 的作用，把 $a_1$ 置为 $0$。于是，式 \textup{(11.14)} 是两个力共同作用的结果：
\begin{equation}
a_1=g+f/m_1, \tag{11.15}
\end{equation}
其中 $f=-gm_1$，这与式 \textup{(11.14)} 完全相同。

同样的操作可以施于 Cartwright 的化油器：例如，在乳化管处装一个流量调节器，就可以在不改变腔体几何形状的情况下把气体流出量固定住。这一操作也绝对适用于经济系统——在那里，大多数方程背后都站着人；方程的左侧可以通过让行为人接触不同信息来固定，而不必改变物理世界中的参数。一个典型例子出现在就业歧视案件中（\S4.5.3）：要检验“性别对录用的效应”，无须在生理上改变申请人的性别，只需改变雇主对申请人性别的知晓即可。我至今尚未见到一个在此处所描述的意义上不具模块性的经济系统。

这种“在方程右侧增加一项以确保左侧恒定”的操作，正是 Haavelmo (1943) 在经济情境中设想手术的方式。他的智慧为何在 2008 年从其弟子的教义中消失，是经济学的一大谜团（见 Hoover (2004)）；我的理论仍然不变（第 138 页）：这一切都源于一次草率的记号选择，而该记号在 1970 年代初统计思想的凌厉入侵下崩解了。

关于计量经济学中的混乱，以及当代计量经济学家为何不愿恢复 Haavelmo 的智慧，将在 \S11.5.4 讨论。

\section{线性结构模型中的因果分析}

\subsection{参数识别的一般准则（第五章，第 149—154 页）}

\qhead{问题（致作者）}

\S5.3.1 所描述的参数识别方法，依赖于对两条基本准则的反复施用：(1) 定理 5.3.1 的单门准则；(2) 定理 5.3.2 的后门准则。该方法在把图中各段的结果组合起来时，可能需要相当繁重的簿记工作。是否存在单一的图形识别准则，能统一这两条定理，从而免去大部分簿记？

\qhead{作者答复}

下述引理给出了一个统一的准则（Pearl 2004）：

\begin{lemma}[直接效应的图形识别]\label{lem:11.5.1}
设 $c$ 表示因果图 $G$ 中赋予箭头 $X\to Y$ 的路径系数。若存在一对 $(W,Z)$，其中 $W$ 是 $G$ 中的单个节点（不排除 $W=X$），$Z$ 是 $G$ 中的一个（可为空的）节点集，使得：
\begin{enumerate}[label=\arabic*.,leftmargin=1.8em]
\item $Z$ 由 $Y$ 的非后代组成；
\item 在从 $G$ 中删去 $X\to Y$ 所得的图 $G_c$ 中，$Z$ 将 $W$ 与 $Y$ d-分离；
\item 在 $G_c$ 中给定 $Z$ 时，$W$ 与 $X$ 是 d-连通的；
\end{enumerate}
则参数 $c$ 可识别。此外，由 $(W,Z)$ 诱导的估计量为
\[
c=\frac{\operatorname{cov}(Y,W\mid Z)}{\operatorname{cov}(X,W\mid Z)}.
\]
\end{lemma}

其直觉是：在给定 $Z$ 的条件下，$W$ 相对于 $X\to Y$ 起着工具变量的作用。\footnote{另见 McDonald (2002a)。更一般的识别方法见 Brito and Pearl (2002a,b,c; 2006)，综述见 Brito (2010)。}

\subsection{结构系数的因果解释}

\qhead{问题（致作者）}

针对 \S5.1 与 \S5.4 所作的断言——正确的因果解释在 SEM 的书籍与论文中显著缺席，包括 1970—1999 年间经济学方面的全部教材——两位读者来信说，“单位变化”的解释在 SEM 文献中其实很常见，也被广泛接受。阿尔伯塔大学的 L.H. 写道：

\begin{quote}\kaishu
L. Hayduk《Structural Equation Modeling with LISREL: Essentials and Advances》(1987) 第 245 页[指出]，斜率可以这样解释：在方程中其他变量保持其原值不变的情况下，$x$ 每变化一个单位时，预计 $y$ 会随之发生的变化量。
\end{quote}

O.D. Duncan 的《Introduction to Structural Equation Models》(1975) 第 1、2 页对 $b$ 的因果含义说得相当清楚。更确切地说，书中写道“$x$ 变化一个单位……使 $y$ 变化 $b$ 个单位”（第 2 页）。我猜 H. M. Blalock 的《Causal models in the social sciences》与 D. Heise 的《Causal analysis》多半也把 $b$ 说成因果的。

佐治亚理工学院的 S.M. 表示同意：

\begin{quote}\kaishu
《Causal Analysis》(1975) 的作者 Heise 认为，因果方程中的 $b$ 就是原因变量变化一个单位时在效应变量上产生的效应量。这是一个被广泛接受的想法。
\end{quote}

\qhead{作者答复}

“单位变化”的想法确实零星地、非正式地出现在若干 SEM 出版物中，然而，与计量经济学中的反事实（\S11.5.5）类似，它从未通过精确的数学定义被操作化。

上面引用的那段文字（出自 Hayduk 1987）可以说明：单位变化的观念在 SEM 文献中通常是如何引入的，以及用现代对因果建模的理解又应当如何引入。原文写道：

\begin{quote}\kaishu
把结构系数解释为“效应系数”，起源于诸如
\[
X_0=a+b_1X_1+b_2X_2+b_3X_3+e
\]
这样的普通回归方程，其中 $X_1$、$X_2$、$X_3$ 是作用于变量 $X_0$ 的变量。我们可以把 $b_1$ 的估计解释为：在 $X_2$ 与 $X_3$ 保持其原值不变的情况下，$X_1$ 每增加一个单位时，预计 $X_0$ 会随之发生的变化量。我们避免以“保持恒定”一语收尾，因为对于含多个方程的模型，这个说法必须放弃，下文将会看到。对 $b_2$ 与 $b_3$ 也应有平行的解释。（Hayduk 1987, p.~245）
\end{quote}

这段话表明，SEM 文献常常混淆两个基本区分。第一个是结构系数与这些系数的回归（或统计）估计之间的区分。我们很少看到前者被独立于后者加以定义——这种混淆十分普遍，在计量经济学教材中尤为难堪。第二个是“保持恒定”与“任其不变”或“发现其保持不变”之间的区分，$\dox{x}$ 算子正是为此而设。（相关的区分还有：“做”与“看”，“干预性变化”与“自然变化”。）

为强调这些区分的中心地位，下面提出对 Hayduk 那段文字的一处简明修订：

\qhead{建议的修订稿}

把结构系数解释为“效应系数”，与诸如
\begin{equation}
X_0=a+b_1X_1+b_2X_2+b_3X_3+e \tag{11.16}
\end{equation}
这样的回归方程中系数的解释有某些相似之处，但有根本区别。

若 \textup{(11.16)} 是回归方程，则 $b_1$ 表示：在 $X_2$ 与 $X_3$ 保持其原值不变的那些情形下，$X_1$ 每变化一个单位时，预计 $X_0$ 会随之发生的变化。我们用条件期望形式地表达这一解释：
\begin{equation}
b_1=E(X_0\mid x_1+1,x_2,x_3)-E(X_0\mid x_1,x_2,x_3)=R_{X_0X_1\cdot X_2X_3}. \tag{11.17}
\end{equation}
注意，作为回归方程，\textup{(11.16)} 是无断言的；也就是说，它不能被任何实验否证，而且由 \textup{(11.17)}，$e$ 自动与 $X_1$、$X_2$、$X_3$ 不相关。

相反，若式 \textup{(11.16)} 表示一个结构方程，它就关于世界作出了经验断言（例如：一旦我们固定 $X_1$、$X_2$、$X_3$，系统中其他变量便不再影响 $X_0$），而 $b_1$ 的解释必须从根本上作两处修改。第一，“$X_1$ 变化一个单位”必须限定为“$X_1$ 发生一个单位的干预性变化”，从而排除由模型中其他变量（可能与 $e$ 相关）所引起的 $X_1$ 变化。第二，“在 $X_2$ 与 $X_3$ 保持不变之处”必须放弃，代之以“若我们使 $X_2$ 与 $X_3$ 保持恒定”，从而即便 $X_2$ 受 $X_1$ 影响也能确保恒定。

形式上，这两处修改表达为：
\begin{equation}
b_1=E\bigl(X_0\mid\dox{x_1+1,x_2,x_3}\bigr)-E\bigl(X_0\mid\dox{x_1,x_2,x_3}\bigr). \tag{11.18}
\end{equation}

“任其保持原值”一语可能导致歧义。让变量任其不变，就允许这些变量发生变化（例如响应 $X_1$ 的单位增加或其他影响），此时 $X_0$ 的变化将对应于总效应
\begin{equation}
E\bigl(X_0\mid\dox{x_1+1}\bigr)-E\bigl(X_0\mid\dox{x_1}\bigr) \tag{11.19}
\end{equation}
或对应于边缘条件期望
\begin{equation}
E(X_0\mid x_1+1)-E(X_0\mid x_1), \tag{11.20}
\end{equation}
取决于 $X_1$ 的变化是干预性的还是观测性的。无论我们把 \textup{(11.16)} 当作结构方程还是回归方程，\textup{(11.19)} 与 \textup{(11.20)} 都不符合 $b_1$ 在式 \textup{(11.16)} 中的含义。

\textup{(11.18)} 所表达的解释在所有模型中都成立，包括含多个方程的模型、递归与非递归模型，无论 $e$ 是否与模型中其他变量相关，也无论 $X_2$ 与 $X_3$ 是否受 $X_1$ 影响。相反，表达式 \textup{(11.17)} 只在非常特殊的条件下（由定理 5.3.1 的单门准则所界定）才与 \textup{(11.18)} 一致。正因如此，我们认为 $b_1$ 的“解释”是 \textup{(11.18)} 而非 \textup{(11.17)}：\textup{(11.17)} 解释的是 $b_1$ 的“回归估计”（它很可能有偏），而 \textup{(11.18)} 解释的是 $b_1$ 本身。

\subsection{为 SEM 的因果解释辩护（或曰 SEM 求生指南）}

\qhead{问题（致作者）}

英格兰萨里郡的 J. Wilson 问道：在那些不认可结构方程模型（SEM）之因果解释的考官面前，该如何为他的博士论文辩护。他抱怨说：“博士培养方案完全不重视如何在答辩中为因果解释与政策含义辩护……倘若在培养方案一开始就把因果性讲透，我那篇论文的七万字本可以字字斟酌，先为因果假设辩护，再为数据辩护，最后为解释辩护……（不知道这个问题有多普遍？）回到眼前这项紧迫任务——设法让考官满意，尤其是那两位极难缠的统计学教授，他们似乎在比谁更刻薄。”

\qhead{作者答复}

你所抱怨的现象，正是我写第五章的起因——SEM 的因果解释对大多数 SEM 研究者、领军人物、教育者与从业者来说仍是一个谜。我在 SEMNET 讨论列表上花了大量时间试图纠正当前的忽视，但只有像你这样的学生才能扭转局面，帮助因果解释重归其在 SEM 研究中的核心地位。

至于你的具体问题——如何在敌视这种解释的刻薄考官面前为 SEM 的因果解释辩护——请允许我勾画一个假想场景，其中的你要在一位敌意考官 EX 博士面前为论文的因果解释辩护。（与 EX 博士的任何相似纯属巧合。）

\begin{center}
{\heiti 与敌意考官的对话}\\[0.2em]
{\heiti 或曰：SEM 求生指南}
\end{center}

为简单起见，假定你论文中的模型只含两个方程，
\begin{equation}
y=bx+e_1, \tag{11.21}
\end{equation}
\begin{equation}
z=cy+e_2, \tag{11.22}
\end{equation}
其中 $e_2$ 与 $x$ 不相关。相应的图见图 11.14。再假定你论文的目标是估计参数 $c$，你用最好的 SEM 方法令人满意地估计出 $c=0.78$，并且你对你的发现给出了因果解释。

这时，你那位刻薄的考官 EX 博士登场，质疑你的解释。

\said{EX 博士：}你说的“$c$ 具有因果解释”是什么意思？

\said{你：}我的意思是：$y$ 变化一个单位，将使 $E(Z)$ 变化 $c$ 个单位。

\said{EX 博士：}“变化”和“使……发生”这些词让我不舒服；咱们讲点科学。你的意思是 $E(Z\mid y)=cy+a$ 吗？？？后面这个表达式我能懂，因为给定 $y$ 时 $Z$ 的条件期望 $E(Z\mid y)$ 在数学上有良好定义，我也知道如何从数据中估计它。可“变化”和“使……发生”对我来说是行话。

\said{你：}我说的确实是“变化”，不是“条件期望的增加”；而所谓“变化”，我指的是：如果我们拥有把 $y$ 固定于某个常数 $y_1$、并把这个常数从 $y_1$ 改变为 $y_2$ 的物理手段，那么观测到的 $E(Z)$ 变化将是 $c(y_2-y_1)$。

\said{EX 博士：}哎呀呀，我们是不是有点形而上学了？我在统计学课里从没听说过“固定”。

\said{你：}哦，抱歉，我没意识到您有统计学背景。那样的话，让我把解释稍作改写如下：如果我们有办法做一个受控的随机化实验，把 $y$ 随机化，那么把对照组设为 $y_1$、实验组设为 $y_2$，观测到的 $E(Z)$ 之差将是 $E(Z_2)-E(Z_1)=c(y_2-y_1)$，无论我们选取什么样的 $y_1$ 与 $y_2$。（$Z_1$ 与 $Z_2$ 分别是对照组与实验组中 $z$ 的测量。）\footnote{万一 EX 博士追问：“这就是全部断言吗？”你应当补充：此外我还断言，随机变量 $Z_1-cy_1$ 的分布与变量 $Z_2-cy_2$ 的分布相同。}

\said{EX 博士：}这听起来接近我能理解的东西了。但你似乎在做一个巨大的跳跃，这让我不安。你的数据是非实验的，而你整个研究里一次实验也没做过。你是在告诉我们，你那套 SEM 练习可以拿观测研究的数据，做一番 LISREL 分析，然后得出一个关于受控随机实验结果将会如何的预测？你在开玩笑吧！！你知道如果我们能用 SEM 魔法取代实验研究，全国能省多少钱吗？

\said{你：}这不是魔法，EX 博士，这是朴素的逻辑。我 LISREL 分析的输入不止是非实验数据，它由两部分组成：(1) 数据；(2) 因果假设。我的结论从这两者逻辑地推出。第二个成分在标准实验研究中是缺席的，而正是它使实验研究如此昂贵。

\said{EX 博士：}什么样的假设？“因果”假设？我从没听说过这种怪物。你能像我们通常表达假设那样，用数学方式把它们表达出来吗——比如说，表达为对联合密度的条件，或协方差矩阵的性质？

\said{你：}因果假设是另一类东西；它们无法用密度函数或协方差矩阵的词汇写出。它们表达在我的因果模型里。

\said{EX 博士：}看着你的模型，式 \textup{(11.21)}—\textup{(11.22)}，我看不到任何新词汇；我看到的全是方程。

\said{你：}这些不是普通的代数方程，EX 博士。这些是“结构方程”。若我们正确地解读它们，它们传达了一组您其实熟悉的假设，即关于在该人群上进行的假想随机化实验之结果的假设——我们称之为“因果”假设或“建模”假设，实属无可奈何，但它们可以被理解为关于该人群在各种随机化实验下之行为的假设。

\said{EX 博士：}等等！既然我开始明白你的因果假设是什么了，我反而比先前更困惑了。如果你允许自己对随机化实验下人群的行为作假设，那何必费劲去做研究？何不直接假设：在把 $y$ 随机化的实验中，观测到的 $E(Z)$ 之差应为 $c'(y_2-y_1)$，其中 $c'$ 随便取个方便的数字，从而省下那几个月痛苦的数据收集与分析。相信你其他未经检验的假设的人，也应当相信你的 $c'(y_2-y_1)$ 假设。

\said{你：}并非如此，EX 博士。我的研究纲领开头的那些建模假设，比我的研究结尾处那个断言 $E(Z_2)-E(Z_1)=0.78(y_2-y_1)$ 要温和得多。第一，我的建模假设是定性的，而我的结论是定量的，承诺了 $c=0.78$ 这个具体数值。第二，许多研究者（包括您，EX 博士）会愿意接受我的假设而不接受我的结论，因为前者符合常识理解以及关于世界如何运转的一般理论知识。第三，我的大多数假设都可以用不涉及 $y$ 随机化的实验来检验。这意味着，如果随机化 $y$ 代价高昂或不可行，我们仍可通过控制其他不那么难办的变量来检验这些假设。最后——尽管本研究并非如此——建模假设往往带有一些可在非实验研究中检验的统计蕴含；若检验成功（我们称之为“拟合”），就进一步确认了这些假设的有效性。

\said{EX 博士：}这倒有意思了。让我看看那些“因果”假设或建模假设，好判断它们有多温和。

\said{你：}容易，请看我们的模型，图 11.14。其中

\begin{center}
\begin{tabular}{r@{\quad}l}
$z$ & ——学生期末考试的分数，\\
$y$ & ——学生用于家庭作业的小时数，\\
$x$ & ——家庭作业在期末成绩中的权重（由教师公布）。
\end{tabular}
\end{center}

当我把这个模型写在纸上时，心里想的是两个随机化实验：一个把 $x$ 随机化（即教师随机分配权重），另一个把实际用于家庭作业的时间（$y$）随机化。我在设想这些实验时所立的假设是：

\begin{enumerate}[label=\arabic*.,leftmargin=1.8em]
\item $y$ 的线性性与排除性：
\[
E(Y_2)-E(Y_1)=b(x_2-x_1),
\]
其中 $b$ 未知（$Y_2$ 与 $Y_1$ 分别是在公布权重为 $x_2$ 与 $x_1$ 时会花在家庭作业上的时间）。此外，通过把 $z$ 从方程中排除，我假设分数 $z$ 不影响 $y$，因为在决定 $y$ 的时候 $z$ 尚不可知。
\item $z$ 的线性性与排除性：
\[
E(Z_2)-E(Z_1)=c(y_2-y_1)
\]
对一切 $x$ 成立，其中 $c$ 未知。用文字说：除经由 $y$ 之外，$x$ 对 $z$ 没有影响。
\end{enumerate}

此外，我还对非实验条件下支配 $x$ 的未测因素作了定性假设：我假设 $x$ 与 $z$ 没有共同原因。EX 博士，您对这些假设有异议吗？

\said{EX 博士：}好吧，我同意这些假设比你论文结论 $E(Z_2)-E(Z_1)=0.78(y_2-y_1)$ 那句生硬而缺乏支持的宣示要温和，而且我有点惊讶：如此温和的假设竟能支撑起一个关于家庭作业对分数之实际效应的、大胆的预测（在实验设定下）。但我对你的共同原因假设仍不满意。在我看来，一位强调家庭作业重要性的教师，多半也是一位鼓舞人心、卓有成效的教师，所以 $e_2$（其中包含教学质量等因素）应当与 $x$ 相关，这与你的假设相反。

\said{你：}EX 博士，您现在开始像一位 SEM 研究者那样说话了。我们不再攻击方法及其哲学，而开始讨论实质问题——例如，“教师的成效与其布置家庭作业的权重不相关”这一假设是否合理。我个人就遇到过对家庭作业毫不在意的优秀教师，反之亦然。

但这不是我论文的主题。我并不主张教师的成效与其给家庭作业的权重不相关；我把这留给其他研究者在未来研究中去检验（或许已经被检验过了？）。我所主张的仅仅是：那些愿意接受“教师成效与其家庭作业权重不相关”这一假设的研究者，将会发现下面这一点很有意思——该假设与数据相结合，逻辑上蕴含这一结论：每天多做一小时家庭作业，会使学生分数（平均）提高 0.78 分。而如果我们被允许做一个把家庭作业量（$y$）随机化的受控实验，这一主张便可以得到经验验证。

\said{EX 博士：}我很高兴你并不坚持你的建模假设为真；你只是陈述它们的似真性并阐明其推论。这一点我无法反对。但我还有一个问题。你说你的模型没有任何统计蕴含，因此无法检验其与数据的拟合程度。你怎么知道的？这难道不让你不安吗？

\said{你：}我只要看看图、检查缺失的连线就知道了。有一条名叫 d-分离的准则（见 \S11.1.2，“无泪的 d-分离”）使 SEM 的学生能够瞥一眼图就判定：相应模型是否蕴含某种以变量间偏相关消失为形式的约束。大多数（虽非全部）统计蕴含都属于这种性质。我们例子中的模型对协方差矩阵不蕴含任何约束，因此它可以完美地拟合任何数据。我们把这种模型称为“饱和的”。有些 SEM 研究者甩不掉统计检验的传统，把这一特征视为模型的缺陷。其实不然。手头有一个饱和模型，只说明研究者不愿作出不可信的因果假设，而他所愿作出的那些温和假设又太弱，产生不出统计蕴含。这样一种审慎态度应当受到称许，而不是谴责。诚然，如果我的模型不饱和——比如说，如果 $e_1$ 与 $e_2$ 不相关——我会高兴。但本情形并非如此；常识告诉我们 $e_1$ 与 $e_2$ 相关，数据也显示了这一点。我曾试着假设 $\operatorname{cov}(e_1,e_2)=0$，结果拟合极差。难道我要为了给模型封个“非饱和”的爵位而作出无根据的假设吗？不！我宁愿作出合理的假设，得出有用的结论，并把结果与假设并列报告。

\said{EX 博士：}但假定还有另一个饱和模型，建立在同样可信的假设之上，却导出一个不同的 $c$ 值。你难道不该担心：你的一些初始假设可能是错的，因而你的结论 $c=0.78$ 也是错的？数据中没有任何东西能帮你偏好其中一个模型。

\said{你：}我确实担心，而且事实上，我可以立刻枚举出所有这类竞争模型的结构；图 11.15 中的两个模型就是例子，此外还有很多。（这同样可以用 d-分离准则来做；见第 145—148 页。）但请注意，竞争模型的存在丝毫不会削弱我先前的断言：“接受模型 $M$ 之定性假设的研究者，不得不接受结论 $c=0.78$。”这一断言在逻辑上依然无懈可击。更何况，这一断言还可以进一步细化：把每个竞争模型的结论连同其背后的假设一并报告。结论的格式于是变成：

\begin{center}
\begin{tabular}{ll}
若你接受假设集 $A_1$，则蕴含 $c=c_1$，\\
若你接受假设集 $A_2$，则蕴含 $c=c_2$，\\
……以此类推。
\end{tabular}
\end{center}

\said{EX 博士：}我明白了。但仍有一个问题：万一我们希望超越这些条件性陈述，真的在各假设集之间作出抉择，难道没有 SEM 方法能助一臂之力吗？我们搞统计的人不习惯面对两个相互竞争的假设却无法对之施加某种检验——哪怕是很弱的检验。

\said{你：}这正是统计分析数据与 SEM 之间的一个根本区别。统计假设按定义就是可由统计方法检验的。相反，SEM 模型建立在因果假设之上，而因果假设同样按定义（见第 39 页）无法接受统计检验。如果我们的两个竞争模型都是饱和的，我们事先就知道除了以条件格式报告结论（如上所列）之外别无他法。然而，如果竞争发生在几个同样可信但统计上可区分的模型之间，我们就面对着一个百年之久的老问题——模型选择，其中已有 AIC 等各种选择准则被提出来用于分析。不过，模型选择问题如今被赋予了一个新的因果转向：我们的使命不是最大化拟合度，也不是最大化预测力，而是对 $c$ 这类因果参数给出最可靠的估计。这完全是一个新的竞技场（见 Pearl 2004）。

\said{EX 博士：}有意思。现在我明白为什么我的统计学家同行们一遇到 SEM 方法就如此彻底地困惑、不信任、甚至敌视了（例如 Freedman 1987; Holland 1988; Wermuth 1992）。最后一个问题。你是在发现我是统计学家之后才开始谈随机化实验的。那么你会如何向一个非统计学家解释你的 SEM 策略？

\said{你：}我会用大白话说：“如果我们拥有把 $y$ 固定于某个常数 $y_1$、并把这个常数从 $y_1$ 改变为 $y_2$ 的物理手段，那么观测到的 $E(Z)$ 变化将是 $c(y_2-y_1)$。”多数人懂得“固定”是什么意思，因为这正是政策制定者心里想的。例如，一位关心家庭作业对成绩之效应的教师，并不会从随机化家庭作业的角度去想问题。随机化只是预测固定之效应的一种间接手段。

其实，如果与我交谈的人真正开明（许多统计学家确实如此），我甚至会动用反事实词汇，比如说：一个在花了 $y$ 小时做家庭作业之后考得 $z$ 分的学生，如果他/她花了 $y+1$ 小时，本会考得 $z+c$ 分。说实话，我在写下方程 $z=cy+e_2$ 时，心里真正想的就是这个；其中 $e_2$ 代表该学生一切在我们的模型中没有被赋予变量名、且不受 $y$ 影响的其他特征。我甚至没有去想 $E(Z)$，只想着一个典型学生的 $z$。反事实是我们拥有的、用以表达科学关系之含义的最精确的语言工具。但我在与统计学家交谈时会避免提及反事实，因为——这很遗憾——统计学家往往怀疑确定性概念，或怀疑那些无法立即检验的概念，而反事实正是这类概念（Dawid 2000; Pearl 2000）。

\said{EX 博士：}谢谢你让我了解了 SEM 的这些方面。没有其他问题了。

\said{你：}这是我的荣幸。

\begin{figure}[htbp]
\centering
\begin{tikzpicture}[causal,scale=1.0,every node/.style={transform shape}]
  \node[circle,fill,inner sep=1.4pt] (x) at (0,0) {};
  \node[font=\small,left=0.6mm] at (0,0) {$x$};
  \node[circle,fill,inner sep=1.4pt] (y) at (1.8,0) {};
  \node[font=\small,below=0.6mm] at (1.8,0) {$y$};
  \node[circle,fill,inner sep=1.4pt] (z) at (3.6,0) {};
  \node[font=\small,right=0.6mm] at (3.6,0) {$z$};
  \draw[ed] (x)--(y) node[midway,below=0.8mm,font=\small] {$b$};
  \draw[ed] (y)--(z) node[midway,below=0.8mm,font=\small] {$c$};
  \draw[bid,dashed] (x) to[bend left=48] (y);
  \draw[bid,dashed] (y) to[bend left=48] (z);
\end{tikzpicture}
\caption{式 \textup{(11.21)}—\textup{(11.22)} 背后的图。}
\label{fig:11.14}
\end{figure}

\begin{figure}[htbp]
\centering
\begin{tikzpicture}[causal,scale=1.0,every node/.style={transform shape}]
  % 左：加入直接边 x->z
  \begin{scope}
    \node[circle,fill,inner sep=1.4pt] (x1) at (0,0) {};
    \node[font=\small,left=0.6mm] at (0,0) {$x$};
    \node[circle,fill,inner sep=1.4pt] (y1) at (1.7,0) {};
    \node[font=\small,below=0.6mm] at (1.7,0) {$y$};
    \node[circle,fill,inner sep=1.4pt] (z1) at (3.4,0) {};
    \node[font=\small,right=0.6mm] at (3.4,0) {$z$};
    \draw[ed] (x1)--(y1) node[midway,below=0.8mm,font=\small] {$b$};
    \draw[ed] (y1)--(z1) node[midway,below=0.8mm,font=\small] {$c$};
    \draw[ed] (x1) to[bend left=52] (z1);
  \end{scope}
  % 右：加入混杂弧 x<->z
  \begin{scope}[xshift=6.4cm]
    \node[circle,fill,inner sep=1.4pt] (x2) at (0,0) {};
    \node[font=\small,left=0.6mm] at (0,0) {$x$};
    \node[circle,fill,inner sep=1.4pt] (y2) at (1.7,0) {};
    \node[font=\small,below=0.6mm] at (1.7,0) {$y$};
    \node[circle,fill,inner sep=1.4pt] (z2) at (3.4,0) {};
    \node[font=\small,right=0.6mm] at (3.4,0) {$z$};
    \draw[ed] (x2)--(y2) node[midway,below=0.8mm,font=\small] {$b$};
    \draw[ed] (y2)--(z2) node[midway,below=0.8mm,font=\small] {$c$};
    \draw[bid,dashed] (x2) to[bend left=52] (z2);
  \end{scope}
\end{tikzpicture}
\caption{与图 11.14 统计等价的模型。}
\label{fig:11.15}
\end{figure}

\subsection{今日的经济学建模在何处？——与 Heckman 一起追求原因}

本书 \S5.2 痛惜过去三十年间计量经济学对结构方程建模之理解的衰退（另见 Hoover 2003, “Lost Causes”），并把这一衰退归因于一次草率的记号选择——它模糊了代数方程与结构方程之间的本质区别。在一系列文章里（Heckman 2000, 2003, 2005; Heckman and Vytlacil 2007），James Heckman 着手扭转这一观感，把因果建模重新确立为经济学研究的核心关注，并重振经济学作为因果分析活跃前沿的地位。无论从哪个角度看，这都不是一件易事。采纳邻近学科中涌现的概念与技术进展，等于承认计量经济学数十年的疏忽；而拒斥这些进展，则必须为它们找到计量经济学的替代物。Heckman 选择了后一条路，尽管因果建模的多数现代进展恰恰扎根于 Haavelmo (1943)、Marschak (1950) 以及 Strotz and Wold (1960) 等经济学家的思想。

Heckman 纲领中的一步，是拒斥 do 算子以及它所依据的“手术”语义，从而剥夺了经济学家使用本书所发展的反事实结构语义（尤其是第七章）的机会——而这一语义恰恰把传统计量经济学与潜在结局进路统一起来。Heckman 拒斥手术的理由，可概括如下：

\begin{quote}\kaishu
外部（强制）变量的受控变化，是在非递归模型中定义因果效应的关键……Pearl 用他生动的语言，以“关掉一个方程”或施行“手术”来定义因果效应。他隐含地假定：在联立方程组中施行“手术”或关掉一个方程，会唯一地确定某一个结局变量或内生变量（在我例子中即另一个人的消费）。一般而言，并不会。对一个方程施加约束，就是对整个内生变量集合施加限制。一般而言，联立方程组中没有哪个方程能唯一决定任何单个结局变量。关掉一个方程，也可能影响系统中其他方程的参数，从而违反参数稳定性的要求。（Heckman and Vytlacil 2007）
\end{quote}

显然，Heckman 的反对意见与 Cartwright 的相同（\S11.4.6）：

\begin{enumerate}[label=\arabic*.,leftmargin=1.8em]
\item 理想的手术在技术上可能不可行；
\item 经济系统是非模块的。
\end{enumerate}

我们在前面四个小节（11.4.3—11.4.6）中已驳斥了这些反对意见，读者很容易把它们重新施用于解构 Heckman 的论证。不过需要再次强调：与 Cartwright 的情形一样，这些反对意见源自把（反事实的）定义任务与识别任务、实际估计任务混为一谈——这是研究者中常见的一种混淆，而 Heckman (2005) 本人还严厉警告读者要避免它。

这种混淆在 Heckman 的如下忧虑中尤其明显：“关掉一个方程，也可能影响系统中其他方程的参数。”在物理世界中，试图实现“手术”所规定的条件，有时确实会影响其他方程中的参数；而且我们将看到，Heckman 所提议的“外部变化”同样如此。然而，我们在此处理的是符号操作，不是物理操作。我们的任务，是在一个名为“模型”的符号系统中，为“一个变量对另一个变量的因果效应”建立有意义的数学定义。这使我们得以随意摆弄符号，而无需顾及这些操作在技术上的可行性。实施层面的考量不必进入定义的讨论。

\qhead{因果效应的新定义：“外部变化”}

在缺少手术语义的情况下，Heckman 与 Vytlacil（以下简称 HV）着手构筑一个新的因果效应定义，希望它能免于他们在手术程序中发现的那种种缺陷——办法是把定义建立在“外部变化”（external variation）之上，而不是关掉方程。遗憾的是，他们的新定义——其反事实逻辑的基石——并未得到显式的形式阐述：它被降格为一条半形式的脚注（HV, p.~77），即便是一位好奇而勤奋的读者也会觉得难以辨读。下面是我对 HV 定义在适用于多方程与非线性系统时的推演。

给定方程组
\[
Y_i=f_i(Y,X,U),\qquad i=1,2,\dots,n,
\]
其中 $X$ 与 $U$ 分别是可观测与不可观测的外部变量集合，则 $Y_j$ 对 $Y_k$ 的因果效应分四步计算：

\begin{enumerate}[label=\arabic*.,leftmargin=1.8em]
\item 在 $X$ 中任选一个出现在 $f_j$ 中的成员 $X_t$。若不存在，则以失败退出。
\item 若 $X_t$ 也出现在其他方程中，考虑把它从那些方程中排除（例如，若方程为线性，则把其系数置零；或用常数替换 $X_t$）。\footnote{目前不清楚有什么条件（如果有的话）会禁止人们在例 11.5.2 中令 $e=0$，或干脆忽略 $X$ 而在第二个方程中加入一个虚拟变量 $X'$。HV 给人的印象是：判定 $e$ 能否置零需要对所处理的问题有深刻理解；若这确实是他们的意图，那大可不必。}
\item 解出所得方程组的简化型
\begin{equation}
Y_i=g_i(X,U),\qquad i=1,2,\dots,n, \tag{11.23}
\end{equation}
\item $Y_j$ 对 $Y_k$ 的因果效应由偏导数给出：
\begin{equation}
\frac{\partial Y_k}{\partial Y_j}=\frac{\partial g_k}{\partial X_t}\Big/\frac{\partial g_j}{\partial X_t}. \tag{11.24}
\end{equation}
\end{enumerate}

\begin{example}\label{ex:11.5.2}
考虑三元方程组
\[
Y_1=aY_2+cY_3+eX+U_1,\qquad
Y_2=bY_1+X+U_2,\qquad
Y_3=dY_1+U_3.
\]
所求为 $Y_2$ 对 $Y_1$ 的因果效应。

该系统只有一个外部变量 $X$，它出现在前两个方程中。若能令 $e=0$，则 $x$ 只出现在 $Y_2$ 的方程中，我们便可进入“外部变化”程序的第 3 步。修改后模型的简化型给出
\[
\frac{\partial Y_1}{\partial X}=\frac{a}{1-ba-cd},\qquad
\frac{\partial Y_2}{\partial X}=\frac{1-cd}{1-ab-cd},
\]
于是 $Y_1$ 对 $Y_2$ 的因果效应算得为
\[
\frac{\partial Y_1}{\partial Y_2}=\frac{a}{1-cd}.
\]
作为对照，手术程序构造如下修改后的方程组：
\[
Y_1=aY_2+cY_3+eX+U_1,\qquad
Y_2=y_2,\qquad
Y_3=dY_1+U_3,
\]
由此得到 $Y_2$ 对 $Y_1$ 的因果效应为
\[
\frac{\partial Y_1}{\partial y_2}=\frac{a}{1-cd},
\]
与“外部变化”程序所得的表达式完全相同。
\end{example}

这两套程序很可能总是给出相同的结果；若果真如此，那就赋予了“外部变化”定义以有效性与概念清晰性。

\subsection{外部变化与手术之争}

在把自己的定义与手术程序所提供的定义相比较时，HV 写道（p.~79）：“要在一个相互依存、非递归的系统中定义因果性，或识别因果参数，并不需要关掉一个方程或摆弄参数……更基本的想法是：把不同的外部变量从不同的方程中排除出去；这些变量一旦被操纵，分析者便能构造出所需的因果量。”

在这个问题上，我与 HV 意见不同。我认为“手术”是更基本的想法，动机更扎实，也更适合政策评估任务。我还要指出，把定义建立在排除与外部变化之上，有如下缺陷：

\begin{enumerate}[label=\arabic*.,leftmargin=1.8em]
\item 一般而言，“排除”涉及把某个变量从方程中移除，等于“摆弄参数”。因此，它是“手术”的一种形式——对原方程组的修改——从而会招致人们可能对“手术”提出的同样批评。虽说我们在前面几节已驳斥了这类批评，但仍应指出：如果它有一丝一毫的道理，那么这一批评对两种方法同样适用。
\item 完全依赖外部变量来揭示内部因果关系的想法，根植于识别文献（例如“工具变量”研究）——在那里，这类变量充当“自然的实验”。然而，在定义因果效应的语境中，这一限制并不正当，因为“因果效应”意在量化由新的外部操纵所产生的效应，而这类操纵未必显式地出现在模型中，也未必在研究的收集数据阶段起作用。更何况，每一个因果结构方程模型，按其本性，都通过手术提供了一种模拟此类外部操纵的隐含机制。

的确，大多数政策评估任务所关注的，正是对内生变量施加直接控制的新的外部操纵。举例来说：一位制造商在多年让价格随成本浮动（即 $\text{价格}=f(\text{成本})$）之后，正决定是否把某产品的现价翻倍。这一决定相当于从手头的模型（即产生现有数据的那个模型）中移除方程 $\text{价格}=f(\text{成本})$，代之以一个等于新价格的常数。这一移除忠实地模拟了被评估的决定；而试图借助“外部变量”绕开它，则是人为的，且鲜有帮助。

再举一例：考虑那个研究得很透彻的问题（Heckman 1992）——评估终止一项教育项目所造成的影响，该项目依一组资格条件录取学生。项目终止后，方程 $\text{录取}=f(\text{资格})$ 将不再成立，而没有任何外部变量能模拟这一新条件（即 $\text{录取}=0$），除非它真的使方程 $\text{录取}=f(\text{资格})$ 失效（或“被忽略”“被关掉”）。

同样有趣的是：Haavelmo (1943) 用来定义因果效应的方法，在数学上等价于手术，而不是外部变化。手术要求把方程 $Y_j=f_j(Y,X,U)$ 替换为 $Y_j=y_j$；Haavelmo 则写成 $Y_j=f_j(Y,X,U)+x_j$，其中 $X_j$ 的选取使 $Y_j$ 保持恒定，$Y_j=y_j$。于是，由于 $X_j$ 使 $Y_j$ 摆脱了 $f_j(Y,X,U)$ 的任何残余影响，Haavelmo 的方法等价于手术。Heckman 的外部变化方法则把 $Y_j$ 留在 $f_j$ 的影响之下。
\item 基于外部变化的定义有一个显而易见的缺陷：目标方程可能不含任何可观测的外部变量。事实上，在许多情形下，系统中可观测外部变量的集合是空的（例如图 3.5）。此外，基于两个偏导数之比的定义，不易推广到含离散变量的非线性系统。因此，认真接受 Heckman 定义的人，将失去目前可用于无工具变量模型的诸多识别技术（见第三、四章）；更严重的是，他们将甚至无法追问：在这类模型中因果效应是否可识别——对于未定义的量，识别问题毫无意义。
\end{enumerate}

所幸，一旦认识到定义可以建立在纯符号操作之上，我们便获得了解放：我们可以修改 Heckman 的提议，在我们希望的任何方程中加入虚构的外部变量。所加入的变量随后可用于定义因果效应，其方式类似于式 \textup{(11.23)} 与 \textup{(11.24)} 中的步骤（假定变量连续）。这使我们更接近手术，唯一的根本差别是：$Y_j$ 仍留在 $f_j(Y,X,U)$ 的影响之下。

在论证了基于“外部变化”的定义在概念上动机不足之后，我们现在来探讨它们能否处理非因果的方程组。

\qhead{非因果系统中的方程歧义}

若干经济学家（Leroy 2002; Neuberg 2003; Heckman and Vytlacil 2007）批评 do 算子依赖于因果的、即有方向的结构方程——在其中变量与方程一一对应。HV 是这样表达这一批评的：“一般而言，联立方程组中没有哪个方程能唯一决定任何单个结局变量”（Heckman and Vytlacil 2007, p.~79）。

人们可以猜到，Heckman 与 Vytlacil 在此指的是含无方向方程的系统，即这样一类方程：其中的等号并不表示“由……决定”或“由……引起”这种非对称关系，而表示对称的代数相等。在计量经济学中，这类非因果方程通常表达均衡约束或资源约束；它们在方程两边之间强加相等，却不赋予左侧变量以“结局”变量的特殊地位。

无方向方程的存在，给反事实 $Y_x$ 的手术式定义带来了歧义——该定义要求把决定 $X$ 的方程替换为常数方程 $X=x$。若 $X$ 出现在若干个方程中，且 $X$ 在方程中的位置是任意的，那么这些方程中的每一个都同等有资格被 $X=x$ 替换，而 $Y_x$ 的值（即替换后 $Y$ 的解）将是歧义的。

注意，对称等式在结构上有别于有方向的非递归系统（即含反馈的系统，如图 7.4）中的互为因果，因为在后一类系统中，每个变量恰是一个方程的“结局”。不过，对称约束可以被建模为一个动态反馈系统的解，其中均衡几乎瞬时达到（Lauritzen and Richardson 2002; Pearl 2003a）。

Heckman 与 Vytlacil 给人以这样的印象：方程歧义是手术定义的缺陷，而不会困扰基于排除的定义。事实并非如此。在无方向方程的系统中，我们无从知道该把哪个外部变量从哪个方程中排除，才能得到正确的因果效应。

举例来说：考虑 HV 第 75 页所讨论的一个二元非递归系统：
\begin{equation}
Y_1=a_1+c_{12}Y_2+b_{11}X_1+b_{12}X_2+U_1, \tag{11.25}
\end{equation}
\begin{equation}
Y_2=a_2+c_{21}Y_1+b_{21}X_1+b_{22}X_2+U_2. \tag{11.26}
\end{equation}
假设我们把 $Y_1$ 移到 \textup{(11.26)} 的左侧，得到
\begin{equation}
Y_1=[a_2-Y_2+b_{21}X_1+b_{22}X_2+U_2]/c_{21}. \tag{11.27}
\end{equation}
为定义 $Y_1$ 对 $Y_2$ 的因果效应，我们此时可以在“从 \textup{(11.25)} 中排除 $X_2$”与“从 \textup{(11.27)} 中排除 $X_2$”之间作选择。前者给出 $c_{12}$，后者给出 $1/c_{21}$。可见，我们在为手术挑选方程时所面临的歧义，转化为在为排除挑选方程与外部变量时的歧义。

打破这一歧义的方法由 Simon (1953) 提出，在第 226—228 页有讨论。

\qhead{小结——经济学建模重获生机}

通过排除与操纵外部变量来构造因果量这一想法，在识别问题的语境中固然动机扎实，但当涉及基于模型的定义时，它并没有逻辑基础。相反，基于手术的定义享有普遍性、语义清晰性与计算简洁性。

那么，这一切把计量经济学建模置于何地？“外部变量”进路的失败，对经济分析与政策评估而言是核心问题还是枝节问题？

在其近期几乎每一篇文章中，James Heckman 都强调反事实是经济分析的必要组成部分，也是一个世纪以来计量经济学成就的标志。例如，HV 一文的第一段写道：“它们[政策比较]要求经济学家构造反事实。要预测已在某一环境中试行、而现在拟应用于新环境的政策的效应，要预测新政策的效应，都需要反事实。”同样，在其发表于《Sociological Methodology》的文章（2005）中，Heckman 写道：“自 Haavelmo (1943, 1944) 以来，经济学家就认识到需要精确的模型来构造反事实……计量经济学框架明确说明了反事实如何生成、干预如何指派……”

然而，尽管反事实在计量经济分析中被宣称具有核心地位，好奇的读者却很难在过去四十年中找到哪怕一篇文章或一本教材，其中反事实或因果效应得到了形式定义。所需要的，是在一个提法良好、完全设定的经济模型中计算反事实 $Y(x,u)$ 的程序，其中 $X$ 与 $Y$ 是模型中两个任意变量。通过拒斥 Haavelmo 基于手术的 $Y(x,u)$ 定义，Heckman 使计量经济学又陷入了十年的分裂与歧义，两个对立的阵营几乎完全隔绝地各自工作。

在 Neyman–Rubin 模型的潜在结局框架内工作的经济学家，把反事实当作原初的、不可观测的变量，完全脱离结构方程模型所编码的知识（例如 Angrist 2004; Imbens 2004）。即便是那些应用倾向评分技术的人——其有效性完全依赖于“可忽略性”或“无混杂”这一因果假设——也很少知道如何借助结构知识来确认或否证该假设（见 \S11.3.5）。

在结构方程框架内工作的经济学家（例如 Kennedy 2003; Mittelhammer et al. 2000; Intriligator et al. 1996）则忙于估计参数，同时把反事实当作不应烦扰凡人的形而上学幽灵。他们信赖 Heckman 这样的领袖，指望其精确界定他们费力估计的结构参数具有何种政策含义，并把它们与潜在结局阵营的同事所做的工作联系起来。\footnote{值得注意的是，经济学家 Hoover (2008) 那篇综合性综述文章的参考文献列表，与经济学家 Angrist (2004) 和 Imbens (2004) 的列表几乎不相交——这道裂痕在文化上极深。}

手术语义（第 98—102 页）以及它所蕴含的因果理论（第七至十章），为这两个对立而视野狭窄的计量经济学研究流派提供了一个简单而精确的统一——一个进路中的定理蕴含另一个进路中的定理，反之亦然。经济学家若能复活 Haavelmo (1943)、Marschak (1950) 与 Strotz and Wold (1960) 的基本思想，并以本书所呈现的图与反事实的逻辑使之重获生机，将会大有裨益。

为完整起见，我在此用括号记号明确重申反事实与结构方程之间的两条根本联系。

\begin{enumerate}[label=\arabic*.,leftmargin=1.8em]
\item 反事实的结构定义是：
\[
Y_M(x,u)=Y_{M_x}(u).
\]
读作：对任何模型 $M$ 与背景信息 $u$，反事实条件句“若 $X$ 曾为 $x$，则 $Y$”由 $Y$ 在子模型 $M_x$ 中的解给出（$M_x$ 即 $M$ 的切除版本，其中决定 $X$ 的方程被替换为 $X=x$）。
\item 结构方程 $y=f(x,e(u))$ 的经验断言是：
\[
Y(x,z,u)=f(x,e(u)),
\]
对任何与 $X$ 或 $Y$ 不相交的集合 $Z$ 成立。
读作：若 $X$ 与 $Z$ 曾分别为 $x$ 与 $z$，则 $Y$ 会是 $f(x,e(u))$，与 $z$ 无关，也与模型中其他方程无关。
\end{enumerate}

\section{决策与混杂（第六章）}

\subsection{辛普森悖论与决策树}

\qhead{Nimrod Megiddo（IBM Almaden）来信}

“我不认为‘因果性’是解开该悖论的关键（当然这也是个定义问题），也不认为二十年前不存在审视它的工具。我出身于博弈论，觉得对于喜欢画决策树（decision tree）、并把‘决策’节点与‘机会’节点区分开来的人来说，这个问题并不难。我画了两棵这样的树[图 11.16(a) 与 (b)]，我认为它们澄清了在不同情形下的正确决策。”

\qhead{作者答复}

你针对同一批输入表格构造出两棵不同的决策树，这一事实本身就意味着：构造的关键不在数据之中，而在你从数据背后的故事中获得的某种信息。那是什么信息？

决策树分析的文献确实已存在至少五十年，但据我所知，它从未认真处理过上面提出的问题：“我们究竟用什么信息来引导自己建立正确的决策树？”

我们一致认为，只把频率表交给机器人，不足以完成这项任务。但机器人先生（或统计学家）还需要什么？把故事从 $F=$“女性”换成 $F=$“血压”，对人类来说似乎就够了，因为人们以非形式的方式理解性别与血压在事物格局中所扮演的不同角色。我们能否形式地刻画这些角色，使我们的机器人能够构造出正确的决策树？

我的提议是：把一对 $(T,G)$ 交给机器人（或统计学家、或决策树专家），其中 $T$ 是频率表集合，$G$ 是因果图；于是，瞧，机器人就能自动建立正确的决策树。我说悖论的解答在于因果考量，意思正在于此。更进一步还可以主张：“既然 $(T,G)$ 中的信息已经足够，何不干脆跳过决策树的构造，直接从 $(T,G)$ 得出正确答案？”这正是本书第三、四章的要旨。那么，决策分析那丰富的文献能否从 $(T,G)$ 这种更经济的问题编码中获益呢？影响图（influence diagram；Howard and Matheson 1981; Pearl 2005）的引入是朝这个方向迈出的一步；而正如 \S4.1.2 所示，第二步或许已不远。影响图是针对特定的决策变量与机会变量之选择而参数化的，因果图却不针对任何这样的选择；它蕴含了为任意决策变量选择、以及在决策时刻可能观测到的任意机会变量选择，去参数化一切影响图（与决策树）所需的信息。

最近，Dawid (2002) 发展了一种混合表示，把影响图与贝叶斯网络结合起来：在后者的某些变量上附加一个决策节点（如图 3.2），代表一个潜在的干预，它能够在观测（闲置）模式与干预（激活）模式之间切换（如式 \textup{(3.8)}）。

\begin{figure}[htbp]
\centering
\begin{tikzpicture}[causal,scale=0.95,every node/.style={transform shape},font=\footnotesize]
  \node[circle,draw,inner sep=0.6pt,minimum size=4.6mm] (r) at (4.9,3.6) {};
  \node[rectangle,draw,inner sep=0.6pt,minimum size=4.2mm] (d1) at (2.1,2.7) {};
  \node[rectangle,draw,inner sep=0.6pt,minimum size=4.2mm] (d2) at (7.7,2.7) {};
  \draw[ed] (r)--(d1); \draw[ed] (r)--(d2);
  \node[font=\scriptsize,above left=-1mm] at (3.5,3.15) {男 0.5};
  \node[font=\scriptsize,above right=-1mm] at (6.3,3.15) {女 0.5};
  \node[circle,draw,inner sep=0.6pt,minimum size=3.6mm] (a1) at (0.7,1.8) {};
  \node[circle,draw,inner sep=0.6pt,minimum size=3.6mm] (a2) at (3.5,1.8) {};
  \node[circle,draw,inner sep=0.6pt,minimum size=3.6mm] (a3) at (6.3,1.8) {};
  \node[circle,draw,inner sep=0.6pt,minimum size=3.6mm] (a4) at (9.1,1.8) {};
  \draw[ed] (d1)--(a1); \draw[ed] (d1)--(a2);
  \draw[ed] (d2)--(a3); \draw[ed] (d2)--(a4);
  \node[font=\scriptsize,left=1.2mm] at (1.35,2.3) {$C$};
  \node[font=\scriptsize,right=1.2mm] at (2.85,2.3) {$\neg C$};
  \node[font=\scriptsize,left=1.2mm] at (6.95,2.3) {$C$};
  \node[font=\scriptsize,right=1.2mm] at (8.45,2.3) {$\neg C$};
  \foreach \c/\lx/\rx in {a1/0.0/1.4, a2/2.8/4.2, a3/5.6/7.0, a4/8.4/9.8} {
    \node[circle,draw,inner sep=0.6pt,minimum size=3.2mm] (l\c1) at (\lx,0.8) {};
    \node[circle,draw,inner sep=0.6pt,minimum size=3.2mm] (l\c2) at (\rx,0.8) {};
    \draw[ed] (\c)--(l\c1); \draw[ed] (\c)--(l\c2);
    \node[font=\scriptsize,left=0.5mm] at (\lx-0.28,1.2) {$E$};
    \node[font=\scriptsize,right=0.5mm] at (\rx+0.28,1.2) {$\neg E$};
  }
  \node[font=\scriptsize] at (0.0,0.32) {0.3};
  \node[font=\scriptsize] at (1.4,0.32) {0.2};
  \node[font=\scriptsize] at (2.8,0.32) {0.35};
  \node[font=\scriptsize] at (4.2,0.32) {0.15};
  \node[font=\scriptsize] at (5.6,0.32) {0.1};
  \node[font=\scriptsize] at (7.0,0.32) {0.4};
  \node[font=\scriptsize] at (8.4,0.32) {0.15};
  \node[font=\scriptsize] at (9.8,0.32) {0.35};
\end{tikzpicture}

{\footnotesize\textbf{（a）}对应图 11.17(a) 的决策树。给定男性，$\neg C$ 优于 $C$（$0.35>0.3$）；给定女性，$\neg C$ 同样更优（$0.15>0.1$）。无条件地看，无论男性概率 $p$、女性概率 $1-p$ 取何值，$\neg C$ 仍优于 $C$（$0.35p+0.15(1-p)>0.3p+0.1(1-p)$）。}

\vspace{0.6em}

\begin{tikzpicture}[causal,scale=0.95,every node/.style={transform shape},font=\footnotesize]
  \node[rectangle,draw,inner sep=0.6pt,minimum size=4.4mm] (r) at (4.9,4.2) {};
  \node[circle,draw,inner sep=0.6pt,minimum size=4.4mm] (b1) at (2.1,3.3) {};
  \node[circle,draw,inner sep=0.6pt,minimum size=4.4mm] (b2) at (7.7,3.3) {};
  \draw[ed] (r)--(b1); \draw[ed] (r)--(b2);
  \node[font=\scriptsize,left=1.2mm] at (3.3,3.8) {$C$};
  \node[font=\scriptsize,right=1.2mm] at (6.5,3.8) {$\neg C$};
  \node[circle,draw,inner sep=0.6pt,minimum size=3.6mm] (c1) at (0.7,2.3) {};
  \node[circle,draw,inner sep=0.6pt,minimum size=3.6mm] (c2) at (3.5,2.3) {};
  \node[circle,draw,inner sep=0.6pt,minimum size=3.6mm] (c3) at (6.3,2.3) {};
  \node[circle,draw,inner sep=0.6pt,minimum size=3.6mm] (c4) at (9.1,2.3) {};
  \draw[ed] (b1)--(c1); \draw[ed] (b1)--(c2);
  \draw[ed] (b2)--(c3); \draw[ed] (b2)--(c4);
  \node[font=\scriptsize,above left=-1mm] at (1.55,2.90) {$\neg\text{HBP}\;0.75$};
  \node[font=\scriptsize,above right=-1mm] at (2.65,2.90) {HBP\;0.25};
  \node[font=\scriptsize,above left=-1mm] at (7.15,2.90) {$\neg\text{HBP}\;0.25$};
  \node[font=\scriptsize,above right=-1mm] at (8.25,2.90) {HBP\;0.75};
  \foreach \c/\lx/\rx in {c1/0.0/1.4, c2/2.8/4.2, c3/5.6/7.0, c4/8.4/9.8} {
    \node[circle,draw,inner sep=0.6pt,minimum size=3.2mm] (n\c1) at (\lx,1.2) {};
    \node[circle,draw,inner sep=0.6pt,minimum size=3.2mm] (n\c2) at (\rx,1.2) {};
    \draw[ed] (\c)--(n\c1); \draw[ed] (\c)--(n\c2);
    \node[font=\scriptsize,left=0.5mm] at (\lx-0.28,1.62) {$E$};
    \node[font=\scriptsize,right=0.5mm] at (\rx+0.28,1.62) {$\neg E$};
  }
  \node[font=\scriptsize] at (0.0,0.72) {0.45};
  \node[font=\scriptsize] at (1.4,0.72) {0.3};
  \node[font=\scriptsize] at (2.8,0.72) {0.05};
  \node[font=\scriptsize] at (4.2,0.72) {0.2};
  \node[font=\scriptsize] at (5.6,0.72) {0.175};
  \node[font=\scriptsize] at (7.0,0.72) {0.075};
  \node[font=\scriptsize] at (8.4,0.72) {0.225};
  \node[font=\scriptsize] at (9.8,0.72) {0.525};
\end{tikzpicture}

{\footnotesize\textbf{（b）}对应图 11.17(b) 的决策树。该树可按 (c) 所示压缩。}

\vspace{0.6em}

\begin{tikzpicture}[causal,scale=0.95,every node/.style={transform shape},font=\footnotesize]
  \node[rectangle,draw,inner sep=0.6pt,minimum size=4.4mm] (r) at (4.2,3.1) {};
  \node[circle,draw,inner sep=0.6pt,minimum size=4.4mm] (d1) at (1.4,2.1) {};
  \node[circle,draw,inner sep=0.6pt,minimum size=4.4mm] (d2) at (7.0,2.1) {};
  \draw[ed] (r)--(d1); \draw[ed] (r)--(d2);
  \node[font=\scriptsize,left=1.2mm] at (2.7,2.65) {$C$};
  \node[font=\scriptsize,right=1.2mm] at (5.7,2.65) {$\neg C$};
  \node[circle,draw,inner sep=0.6pt,minimum size=3.2mm] (e1) at (0.0,1.0) {};
  \node[circle,draw,inner sep=0.6pt,minimum size=3.2mm] (e2) at (2.8,1.0) {};
  \node[circle,draw,inner sep=0.6pt,minimum size=3.2mm] (e3) at (5.6,1.0) {};
  \node[circle,draw,inner sep=0.6pt,minimum size=3.2mm] (e4) at (8.4,1.0) {};
  \draw[ed] (d1)--(e1); \draw[ed] (d1)--(e2);
  \draw[ed] (d2)--(e3); \draw[ed] (d2)--(e4);
  \node[font=\scriptsize,left=1.0mm] at (0.55,1.62) {$E$};
  \node[font=\scriptsize,right=1.0mm] at (2.25,1.62) {$\neg E$};
  \node[font=\scriptsize,left=1.0mm] at (4.75,1.62) {$E$};
  \node[font=\scriptsize,right=1.0mm] at (6.45,1.62) {$\neg E$};
  \node[font=\scriptsize] at (-0.3,0.62) {0.5};
  \node[font=\scriptsize] at (3.1,0.62) {0.5};
  \node[font=\scriptsize] at (5.3,0.62) {0.4};
  \node[font=\scriptsize] at (8.7,0.62) {0.6};
\end{tikzpicture}

{\footnotesize\textbf{（c）}压缩后的决策树。显然，$C$ 优于 $\neg C$（因为 $0.5>0.4$）。}
\caption{辛普森悖论对应的三棵决策树。}
\label{fig:11.16}
\end{figure}

\subsection{时间顺序信息足以构造决策树吗？}

\qhead{Megiddo 回信}

“‘因果性’一词把一些本不必出现的问题引入进来，诸如决定论、自由意志、因与果等等。真正要紧的是一项规定：在确定结局的过程中，变量 $X$ 的取值先于变量 $Y$ 的取值被确定，且 $Y$ 依赖于 $X$。你喜欢把这种情形称为因果关系。当然，在数学理论中你可以随便给它起什么名字，但这样一来，人们起初会被诱导去发展某种直觉，而这种直觉可能因你所选名字的外部含义而错了。这种直觉在数学模型之外的解读，往往具有现实含义，而这些含义可能是错的——例如，‘$X$ 确实导致 $Y$’。决策树是展示额外的时间顺序信息的一种简单方式，而简单的有向图当然可以更简洁地编码该信息。当你必须精确界定这些图意味着什么时，你会求助于像树那样更完整的描述。总之，我唯一的异议是针对‘因果性’一词的用法；我从未怀疑过，时间顺序信息对于基于过去数据作出正确决策至关重要。”

\qhead{作者答复}

\begin{enumerate}[label=\arabic*.,leftmargin=1.8em]
\item 在 21 世纪之前，把“因果”这类标签加在某些数学关系上确有一定危险——这类标签满载着直觉、神秘与争议。如今情况已非如此：神秘感已消散，而这种解释的现实含义并不错——$X$ 确实导致 $Y$。更何况，如果在把数学搞对之外，人们还有兴趣阐明那些宝贵的直觉，使我们能更精确地解读它们、甚至把它们教给机器人，那就别无选择，只能给这些关系冠以它们在当前语言中已有的名字，即“因果”。
\item 进入决策树的东西，不止时间顺序信息与依赖信息。例如，图 11.17(c) 所传达的时间顺序与依赖信息，与图 11.17(a) 的相同（假定 $F$ 先于 $C$ 发生），然而 (c) 要求一棵不同的决策树（并给出不同的结论），因为在 (a) 中 $F$ 与 $Y$ 之间的依赖是“因果的”，而在 (c) 中是关联性的。因此，若要构造正确的决策树、并给其分支赋予正确的概率，因果考量必须补充时间顺序信息与依赖信息。

作为一项思想实验，设想我们想写一个程序，从图 11.17(a)—(b)—(c) 那样的故事中自动构造决策树。该程序获得经验频率表，并被允许向我们询问 $C$、$E$、$F$ 之间的时间顺序与依赖关系，但不许使用任何因果词汇。这个程序能区分 (a) 与 (c) 吗？答案是：不能。如果无视因果考量，只像你所建议的那样关注时间顺序与依赖信息，那么为图 11.17(c) 构造的决策树，就会与为 $F=$ 性别所画的那棵（图 11.16(a)）完全相同，随之而来的是错误结论：该药（$C$）对 $F=$ 真与 $F=$ 假的患者都有害。这是错的，因为正确答案是：该药对整个总体有益（如血压一例），因此它必定对 $F=$ 真或 $F=$ 假的患者（或二者）有益。决策树掩盖了这一信息，给出了错误的结果。

错误源于给决策树的分支赋予了错误的概率。例如，树中最左的分支被赋予概率 $P(E\mid C,\neg F)=0.6$，这是错的；正确的概率应是 $P\bigl(E\mid\dox{C},\neg F\bigr)$，而在图 11.17(c) 的情形下，它无法由图与表确定。不过我们确实知道 $P\bigl(E\mid\dox{C}\bigr)=0.5$ 与 $P\bigl(E\mid\dox{\neg C}\bigr)=0.4$，因此式 \textup{(6.4)} 或式 \textup{(6.5)} 必有一个被违反。

决策分析中的常规做法，是给树的分支赋予 $P(E\mid \text{行动},Z)$ 这类条件概率，其中 $Z$ 是决策时刻可用的信息。实践者应当被警告不要盲目遵循这一常规。正确的赋值当然是 $P\bigl(E\mid\dox{\text{行动}},Z\bigr)$——它或者可由实验研究估计，或者在因果图使该量可识别时由观测数据估计。
\item 我不同意：为了精确界定因果图的含义，我们必须“求助于像树那样更完整的描述”。在 \S7.1 中我们可以找到把因果图定义为函数集合的形式定义，而这个定义并未援引任何决策树（至少没有显式援引）。因此，因果图有其自身的含义，独立于它可能帮助我们构造的众多决策树。以此类推，说微分方程（比如粒子运动的方程）的含义在于服从该方程的轨迹集合，会显得别扭（尽管在数学上并不错）；微分方程中每一项的含义都可以从基本原理出发被精确界定，而无须对它的解有丝毫了解。
\end{enumerate}

\begin{figure}[htbp]
\centering
\begin{tikzpicture}[causal,scale=0.86,every node/.style={transform shape},font=\footnotesize]
  % (a) F -> C, F -> E, C -> E
  \begin{scope}
    \node[circle,fill,inner sep=1.2pt] (aC) at (0.2,0.9) {};
    \node[font=\scriptsize,left=0.4mm] at (0.2,0.9) {$C$};
    \node[circle,fill,inner sep=1.2pt] (aF) at (1.8,0.9) {};
    \node[font=\scriptsize,right=0.4mm] at (1.8,0.9) {$F$};
    \node[circle,fill,inner sep=1.2pt] (aE) at (1.0,-0.7) {};
    \node[font=\scriptsize,below=0.4mm] at (1.0,-0.7) {$E$};
    \node[font=\scriptsize,above=0.2mm] at (0.2,1.35) {治疗};
    \node[font=\scriptsize,right=0.2mm] at (2.05,0.55) {性别};
    \node[font=\scriptsize,left=0.2mm] at (0.35,-0.75) {康复};
    \draw[ed] (aF)--(aC); \draw[ed] (aC)--(aE); \draw[ed] (aF)--(aE);
    \node[font=\scriptsize] at (1.0,-1.25) {(a)};
  \end{scope}
  % (b) C -> F, C -> E, F -> E
  \begin{scope}[xshift=4.6cm]
    \node[circle,fill,inner sep=1.2pt] (bC) at (0.2,0.9) {};
    \node[font=\scriptsize,left=0.4mm] at (0.2,0.9) {$C$};
    \node[circle,fill,inner sep=1.2pt] (bF) at (1.8,0.9) {};
    \node[font=\scriptsize,right=0.4mm] at (1.8,0.9) {$F$};
    \node[circle,fill,inner sep=1.2pt] (bE) at (1.0,-0.7) {};
    \node[font=\scriptsize,below=0.4mm] at (1.0,-0.7) {$E$};
    \node[font=\scriptsize,above=0.2mm] at (0.2,1.35) {治疗};
    \node[font=\scriptsize,right=0.2mm] at (2.0,0.6) {血压};
    \node[font=\scriptsize,left=0.2mm] at (0.35,-0.75) {康复};
    \draw[ed] (bC)--(bF); \draw[ed] (bC)--(bE); \draw[ed] (bF)--(bE);
    \node[font=\scriptsize] at (1.0,-1.25) {(b)};
  \end{scope}
  % (c) latent U -> C, U -> F, C -> E, F <-> E (associational)
  \begin{scope}[xshift=9.4cm]
    \node[circle,fill,inner sep=1.2pt] (cC) at (0.2,0.9) {};
    \node[font=\scriptsize,left=0.4mm] at (0.2,0.9) {$C$};
    \node[circle,fill,inner sep=1.2pt] (cF) at (1.8,0.9) {};
    \node[font=\scriptsize,right=0.4mm] at (1.8,0.9) {$F$};
    \node[circle,fill,inner sep=1.2pt] (cE) at (1.0,-0.7) {};
    \node[font=\scriptsize,below=0.4mm] at (1.0,-0.7) {$E$};
    \node[circle,draw,inner sep=0.5pt,minimum size=2.4mm] (cU) at (1.5,2.15) {};
    \node[font=\scriptsize,above=0.2mm] at (0.2,1.35) {治疗};
    \draw[ed] (cC)--(cE);
    \draw[edd] (cU)--(cC); \draw[edd] (cU)--(cF);
    \draw[bid,dashed] (cF) to[bend left=28] (cE);
    \node[font=\scriptsize] at (1.0,-1.25) {(c)};
  \end{scope}
  % (d) C <-> F dashed, C -> E, F -> E
  \begin{scope}[xshift=14.2cm]
    \node[circle,fill,inner sep=1.2pt] (dC) at (0.2,0.9) {};
    \node[font=\scriptsize,left=0.4mm] at (0.2,0.9) {$C$};
    \node[circle,fill,inner sep=1.2pt] (dF) at (1.8,0.9) {};
    \node[font=\scriptsize,right=0.4mm] at (1.8,0.9) {$F$};
    \node[circle,fill,inner sep=1.2pt] (dE) at (1.0,-0.7) {};
    \node[font=\scriptsize,below=0.4mm] at (1.0,-0.7) {$E$};
    \draw[ed] (dC)--(dE); \draw[ed] (dF)--(dE);
    \draw[bid,dashed] (dC) to[bend left=40] (dF);
    \node[font=\scriptsize] at (1.0,-1.25) {(d)};
  \end{scope}
\end{tikzpicture}
\caption{说明时间顺序信息不足以构造决策树的图。}
\label{fig:11.17}
\end{figure}

\subsection{Lindley 论因果性、决策树与贝叶斯主义}

\qhead{问题（致作者，来自 Dennis Lindley）}

如果你的假设——把 $X$ 控制在 $x$ 等价于移除 $X$ 的函数并在各处把 $X$ 置为 $x$——成立，那么它是说得通的。我目前不明白的是，这与决策树有何关系。在决策节点上，人们对决策时刻已知的量取条件；在随机节点上，人们在已知条件下纳入一切相关的不确定量。除了联合分布（以及效用考量），别无所需。例如在医学情形中，混杂因子在作出治疗决定时可能已知、也可能未知，而这决定了树的结构。因果性可能介入之处，是当我们用数据去评估树中所需的概率的时候；正是在这里，Novick 与我使用了可交换性。贝叶斯范式把“作为信念的概率”与“作为频率的概率”作了尖锐区分，称后者为机遇。若我对因果性的理解不错，那么在这种语境下，我们的概念被你的概念便利地取代，应是合理的。

\qhead{作者答复}

许多决策分析者持这样的立场：因果性并不必要，因为“除了联合分布（以及效用考量），别无所需”（见 \S11.6.2 的讨论）。他们既对也不对。概率分布确实就是所需的一切，但所需的分布 $P(y\mid\dox{x},\text{see}(z))$ 属于一种特殊的类型；除非我们用因果知识（例如因果图所提供的知识）来补充数据，否则它们不能从支配我们数据的联合分布 $P(y,x,z)$ 推出。

你的下一句话道出了一切：

\begin{quote}\kaishu
因果性可能介入之处，是当我们用数据去评估树中所需的概率的时候，……
\end{quote}

我认为频率／信念之分与因果性的讨论关系不大。假定辛普森故事中的那些表不是频率表，而是某人对各种联合事件 $(C,E,F)$、$(C,E,\neg F)$……等等发生情况的主观信念的汇总。我的断言依然成立：这一信念汇总不足以构造我们的决策树。我们还需要评估自己对假想事件“若采取决定 $\dox{C}$，$E$ 将会发生”的信念；而如我们所见，仅凭时间信息不足以从表格化的信念汇总中导出这一评估。因此，我们无法从这一信念汇总构造决策树；我们需要一项额外的成分，我称之为“因果”信息，而你选择称之为“可交换性”——我不愿在命名上争执，只想指出：如果我们信任人类能提供可靠的信息，我们就应当尊重他们表达该信息所用的词汇；而令谨慎的学者们大失所望的是，那个词汇明确无误是因果的。

\qhead{问题更新（致作者）}

你关于概率与决策树的论点很有道理，我同意你在这里所说的；这一点我此前未曾领会。谢谢。让我把论证重述一遍，看我们是否一致。处理决策树时，容易看清解决问题需要哪些概率；但不太容易看清这些概率如何被数值地评估。不过，我看不出因果性如何完全解决这个问题。当然，可交换性也不能。

\qhead{作者答复}

我很高兴我们已把问题收窄为几个简单而具体的问题：(1) 如何评估决策树所需的概率；(2) 这些概率从何而来；(3) 这些概率如何被经济地编码；也许还有 (4) 这些概率是否必须服从某些内在一致性规则——尤其当我们构造若干棵决策树、而它们在决策节点与机会节点的选择上各不相同的时候。

我欢迎这一收窄，因为它大大有助于我为“明确区分因果关系与概率关系”所作的论证。总的来说，我发现贝叶斯统计学家是最难被说服承认这一区分之必要性的一类统计学家。为什么？因为古典统计学家时刻警惕那些无法由硬数据证实的假设，而贝叶斯主义者在这方面更为宽容——这是对的。然而，一旦把人类判断认可为合法信息来源，贝叶斯主义者便不再那么细致地追踪该信息的性质与来源。你先前那段话，描述的正是贝叶斯研究者中对条件化的习惯性、不加批判的使用；这种使用偶尔会把不慎者引入死胡同（例如 Rubin (2009)）：

\begin{quote}\kaishu
在决策节点上，人们对决策时刻已知的量取条件；在随机节点上，人们在已知条件下纳入一切相关的不确定量。除了联合分布（以及效用考量），别无所需。
\end{quote}

正如纽科姆悖论教给我们的（见 \S4.1），“在决策节点上，人们对决策时刻已知的量取条件”并不完全正确——至少有些“取条件”必须被替换为“做”。若非如此，一切决策树都会变成笑话：“病人应当避免去看医生，以降低自己病重的概率（Skyrms 1980, p.~130）；工人绝不应当赶着上班，以降低自己睡过头的概率；学生不应当为考试复习，以免这证明他们学业落后；如此等等。简言之，一切补救行动都应当被禁止，以免它们增加‘确实需要补救’的概率。”（第四章，第 108 页）

然而，即便逃出了这个“条件化”陷阱，贝叶斯哲学家仍然看不出：为决策节点发出的分支评估概率，与为机会节点发出的分支评估概率，有什么不同。对贝叶斯主义者而言，两种评估都是概率评估。前者涉及对假想实验的心智模拟，后者涉及对被动观测的心智设想，这在他看来无关紧要，因为贝叶斯主义者全神贯注于“作为信念的概率”与“作为频率的概率”之分。

这种专注使他们不太敏感于这样一个事实：信念有各种各样的来源，而区分关于实验结果的信念与关于被动观测结果的信念至关重要（Pearl 2001a, “Bayesianism and Causality”）。第一，后者是稳健的，前者不是；也就是说，如果我们关于被动观测的主观信念有误，随着观测数增加，误差的影响终将被冲刷掉。关于实验结果的信念则不然：错误的先验将产生永久性偏倚，无论样本量多大。第二，关于实验结果的信念无法用概率演算的语法来表达；因此，为形式地表达这类信念并把它们与数据一致地结合，概率演算必须用因果记号（即图、$\dox{x}$ 或反事实）来充实。

迄今为止，我发现了两种赢得贝叶斯主义者之心的有效方式，一种涉及“经济性”概念（见我在 \S11.6.2 的讨论），另一种涉及“一致性”概念。

给定 $n$ 个所关心的变量，可以设想由这些变量构造出数量庞大的决策树，每一棵对应一种不同的时间顺序选择，以及从这些变量中选取决策节点与机会节点的一种不同方式。于是自然产生一个问题：决策者如何确保对所有这些决策树的概率评估都是可复现的？我们当然不能假定人类把所有这些潜在的决策树显式地储存在头脑中。为保证可复现性，我们必须假定：所有这些评估都派生自关于决策与机遇事件的某种经济的知识表示。

于是，因果关系可以被视为这样一种经济表示，决策树由之构造出来。的确，正如我在给 Megiddo 的信中所写（\S11.6.2），如果我们须教会一个机器人按需、依照可用的知识与信念构造这类决策树，最好的办法是给机器人一对输入 $(G,P)$，其中 $G$ 是因果图，$P$ 是我们对所关心变量的联合分布（若我们是贝叶斯主义者，则是主观分布）。借助这对对象，机器人应能一致地构造出所需的全部决策树——对变量划分为决策节点与机会节点的任意划分皆然——并在各分支上复现相应参数。这是贝叶斯主义者欣赏因果性的一种方式，且不触犯“它不过是联合分布而已……”这一传统立场。

第二种进路涉及“一致性”。一致性是贝叶斯主义者引以为傲的东西，因为 De Finetti、Savage 等人曾费力构造定性公理，以防止概率判断流于任性，并迫使信念服从概率演算。

我们可以请贝叶斯哲学家告诉我们：关于联合概率的判断（比如 $P(x,y)$），是否应在某种意义上与关于基于决策之概率的判断（比如 $P(y\mid\dox{x})$）相一致——后者量化的是从带有两个备择 $\dox{X=x}$ 与 $\dox{X=x'}$ 的决策节点发出的分支。我们接着可以问贝叶斯主义者：这些概率与通常的条件概率 $P(y\mid x)$（即在另一棵决策树中，从机遇事件 $X=x$ 发出的、对结局 $Y=y$ 所评估的概率）之间，是否应存在任何联系。

要说服我们的贝叶斯主义者相信这两种评估不可能完全任意、而必须服从某些一致性限制，应该不太难。例如，不等式
\[
P(y\mid\dox{x})\ge P(y,x)
\]
应对一切事件 $x$ 与 $y$ 成立。\footnote{该不等式由 \textup{(3.52)} 或 \textup{(9.33)} 推出。一致性约束的完整刻画见 Tian, Kang, and Pearl (2006)。例如，对任意三个变量 $X,Y,Z$，一致性要求
\[
P\bigl(y\mid\dox{x,z}\bigr)-P\bigl(y,x\mid\dox{z}\bigr)-P\bigl(y,z\mid\dox{x}\bigr)+P(x,y,z)\ge 0.
\]
若因果图的结构已知，则定义 1.3.1 的条件构成对一切一致性要求的完整刻画。}
而且，只要 $P(y\mid\dox{x})$ 是按第三章的规则从因果网络导出的，这类一致性限制便自动得到满足。

这两个论证应当足以吸引贝叶斯主义者开始从因果演算中获取数学上的收益，同时保持审慎与怀疑；正如《塔木德》所说：

\begin{center}
\kaishu “由利益而生理解”
\end{center}

（这是“mitoch shelo lishma, ba lishma”（《塔木德》，Psahim, 50b）的自由意译。）

贝叶斯主义者终将拥抱因果词汇，我对此毫不怀疑。

\subsection{为什么混杂不是一个统计概念？}

2001 年 6 月，我收到两封关于拙文《健康科学中的因果推断》（Pearl 2001c）的匿名审稿意见。审稿人提出的问题令我震惊，因为它们让我想起某些统计学家思考因果性时那种陈旧的方式，也让我想起前方仍有艰巨的教育工作要做。为有助于这一工作，我把自己的答复收入本章。关于因果—统计之分的相关讨论见 \S11.1。

\qhead{审稿意见摘录}

\said{审稿人 1：}“统计概念与因果概念之间的对比被夸大了。随机化、工具变量等等都有清晰的统计学定义。……[该文主张]‘任何系统的因果分析方法都必须要求新的数学记号’。这是错的：医学科学中有着悠久传统，进行非形式——但系统而成功——的因果推断。”

\said{审稿人 2：}“该文作了许多笼统的评论，都依赖于区分‘统计’概念与‘因果’概念……此外，在被列为因果（因而据该文说，是非统计的）概念的清单中，就有混杂——而它牢固地奠基于标准的频率派统计学。统计学家倾向于说这样的话：‘当 $U$ 与 $X$、以及 $U$ 与 $Y$ 都不独立时，$U$ 是考察处理 $X$ 对结局 $Y$ 之效应时的一个潜在混杂因子。’那么，为什么混杂不是一个统计概念？……若作者要我相信这一点，他至少得举出一个例子说明通常的分析是如何失败的。”

\qhead{作者回应}

审稿人 1 似乎主张一种非形式的因果研究途径（不管那是什么意思；它让我想起 Bernoulli 之前的非形式统计学）；既然如此，他/她来到这个论坛时，就已带着对拙文基本目标与原则的先入之见。让我们让历史来评判我们之间的是非。

本文的目标读者是持审稿人 2 那种看法的读者——他们试图把本文的主张（我承认，偶尔是笼统的）与传统的统计智慧调和起来。对这位审稿人，我有如下评论。

你质疑我所提出的统计概念与因果概念之划界的有用性。让我通过你举的例子——混杂——来演示这种有用性。你写道，“混杂牢固地奠基于标准的频率派统计学”，而统计学家倾向于说：“当 $U$ 与 $X$、以及 $U$ 与 $Y$ 都不独立时，$U$ 是考察处理 $X$ 对结局 $Y$ 之效应时的一个潜在混杂因子。那么，为什么混杂不是一个统计概念？”

我书的第六章不惜篇幅地解释：为什么这个定义在充分性检验与必要性检验两方面都失败了，以及为什么这类定义的一切变体都必将在基本原理上失败。我只举几个例子来演示这一点；更多的讨论见 \S11.3.3。考虑一个同时受 $X$ 与 $Y$ 影响的变量 $U$——比如说，只要 $X$ 与 $Y$ 都达到高水平它就取 $1$。$U$ 满足你的判据，然而 $U$ 并不是考察处理 $X$ 对结局之效应时的混杂因子——事实上，$U$ 在我们的分析中完全可以被安全地忽略。（对任何其作用于 $Y$ 的效应被 $X$ 中介的变量 $U$，例如我文中图 2 的 $Z$，也是如此。）第二个例子：考虑一个位于从 $X$ 到 $Y$ 的“因果通路”上的变量 $U$。这个变量同样满足你的判据，却不是混杂因子——一代代统计学家都被警告过（见 Cox 1958）不要对这类变量作调整。有人可能辩称，你的定义只是混杂的一个必要条件而非充分条件。但这同样不成立。第六章（第 185—186 页）描述了一些例子：其中没有任何变量满足你的定义，而处理 $X$ 对结局 $Y$ 的效应仍然是混杂的。

还可以构造一个例子（图 6.5），其中 $U$ 是混杂因子（即必须对之调整才能消除效应偏倚），而 $U$ 与 $X$ 和 $Y$ 都不相关联。

我并非第一个发现混杂与其各种统计“定义”之间不一致的人。Miettinen、Cook、Robins 与 Greenland 自 1980 年代中期起就在流行病学中争论这一点——却毫无效果。研究者继续把可压缩性与无混杂等同起来，继续对错误的变量作调整（Weinberg 1993）。更何况，“任何重要概念（例如随机化、混杂、工具变量）都必须有一个统计学定义”这一流行观念，在健康科学中根深蒂固，以至于时至今日——我书出版已十五个月——资历最高、意图最纯的人仍在问：“那么，为什么混杂不是一个统计概念？”

我相信，任何纠正这一传统的尝试都必然听起来笼统；而只有一场全面的运动，才能根除这些关于混杂、统计学与因果性的误解。统计教育牢牢掌握在持审稿人 1 那种看法的人手中。而像你这样追求理解的人，很少获得公开的论坛去问：“那么，为什么混杂不是一个统计概念？”

同样的论证适用于“随机化”与“工具变量”这两个概念。（具有讽刺意味的是，审稿人 1 权威地宣称这些概念“有清晰的统计学定义”。我倒想递给他一个二元分布 $f(x,y)$，请他告诉我们 $X$ 是否被随机化。）这些概念的任何定义都必须援引无法用分布来定义的因果词汇——根本无路可逃。

这就使我回到重新强调统计—因果划界的有用性，如 \S1.5 所定义。那些认识到随机化、混杂、工具变量等概念必须建立在因果信息之上的人，会警惕地分离并阐明援引这些概念的研究背后的因果假设。相反，那些无视因果—统计之分的人（例如审稿人 1），则不会寻求这样的阐明，并将继续对“通常的分析是如何失败的”视而不见。

\section{反事实的演算}

\subsection{线性系统中的反事实}

我们知道，一般而言，形如 $P(Y_x=y\mid e)$ 的反事实查询在经验上可能可识别，也可能不可识别，即便在实验研究中亦然。例如，因果概率
\[
P\bigl(Y_{x'}=y'\mid x',y'\bigr)
\]
一般而言既不能由观测研究、也不能由实验研究识别（第 290 页，推论 9.2.12）。本节要问的是：线性性假设是否使反事实断言在经验上更为扎实。答案是肯定的：

\noindent\textbf{断言 A.} 形如 $E(Y_x\mid e)$ 的任何反事实查询，在线性因果模型中都是经验可识别的，其中 $e$ 为任意证据。

\noindent\textbf{断言 B.} 只要因果效应 $T$ 可识别，$E(Y_x\mid e)$ 也可识别。

\noindent\textbf{断言 C.} $E(Y_x\mid e)$ 由
\begin{equation}
E(Y_x\mid e)=E(Y\mid e)+T\,\bigl[x-E(X\mid e)\bigr] \tag{11.28}
\end{equation}
给出，其中 $T$ 是 $X$ 对 $Y$ 的总效应系数，即
\begin{equation}
T=\frac{dE[Y_x]}{dx}=E\bigl(Y\mid\dox{x+1}\bigr)-E\bigl(Y\mid\dox{x}\bigr). \tag{11.29}
\end{equation}

断言 A 并不令人惊讶。它已由 Balke and Pearl (1994b) 一般地建立，其中 $E(Y_x\mid e)$ 被用协方差矩阵与结构系数给出了显式表达式，而后者是经验可识别的。

断言 B 使 $E(Y_x\mid e)$ 在所有因果效应 $E(Y_x)$ 可识别的模型中，也成为观测可识别的。

断言 C 为 $E(Y_x\mid e)$ 提供了一个直觉上极具说服力的解读，读作：给定证据 $e$，为计算 $E(Y_x\mid e)$（即在“$X$ 曾为 $x$、而非其当前值”这一假设下 $Y$ 的期望），先算出以证据 $e$ 为条件的 $Y$ 的最佳估计 $E(Y\mid e)$，再加上当 $X$ 从其当前最佳估计 $E(X\mid e)$ 被迫转变到其假设值 $X=x$ 时，预期在 $Y$ 上出现的变化。最后这一加项不是别的，正是效应系数 $T$ 乘以 $X$ 的预期变化，即 $T[x-E(X\mid e)]$。

\noindent\textbf{注：}式 \textup{(11.28)} 也可以用 $\dox{x}$ 记号写作
\begin{equation}
E(Y_x\mid e)=E(Y\mid e)+E\bigl(Y\mid\dox{x}\bigr)-E\bigl(Y\mid\dox{X=E(X\mid e)}\bigr). \tag{11.30}
\end{equation}

\noindent\textbf{证明}（在 Ilya Shpitser 协助下）

不失一般性，假定我们处理的是零均值模型。由于模型是线性的，我们可以把 $X$ 与 $Y$ 之间的关系写作
\begin{equation}
Y=TX+I+U, \tag{11.31}
\end{equation}
其中 $T$ 是 $X$ 对 $Y$ 的总效应，由 \textup{(11.29)} 给出；$I$ 表示含有模型中其他变量（$X$ 的非后代）的项；$U$ 表示外生变量。

总可以通过递归地把每个以 $X$ 为祖先的右侧变量的函数代入，并把一切含 $X$ 的项归并起来构成 $TX$，从而把决定 $Y$ 的函数化为 \textup{(11.31)} 的形式。显然，$T$ 是由 $X$ 出发、终止于 $Y$ 的各条路径之路径代价的 Wright 规则求和（Wright 1921）。

由 \textup{(11.31)} 可以写出
\begin{equation}
Y_x=Tx+I+U, \tag{11.32}
\end{equation}
因为 $I$ 与 $U$ 不受从 $X$ 到 $x$ 这一假设性改变的影响；此外，
\begin{equation}
E(Y_x\mid e)=Tx+E(I+U\mid e), \tag{11.33}
\end{equation}
因为 $x$ 是常数。

\textup{(11.33)} 中的最后一项，可通过对 \textup{(11.31)} 两边取期望来求值，得到
\begin{equation}
E(I+U\mid e)=E(Y\mid e)-T\,E(X\mid e), \tag{11.34}
\end{equation}
把它代入 \textup{(11.33)}，即得
\begin{equation}
E(Y_x\mid e)=Tx+E(Y\mid e)-T\,E(X\mid e), \tag{11.35}
\end{equation}
从而证明了我们的目标公式 \textup{(11.28)}。\hfill$\blacksquare$

$e$ 的三种特殊情形值得注意：

\begin{example}
$e:\ X=x',\ Y=y'$（因果概率的线性对应物，第九章）。由 \textup{(11.28)} 直接得到
\[
E\bigl(Y_x\mid Y=y',X=x'\bigr)=y'+T(x-x').
\]
这在直觉上极具说服力：$Y$ 的假设性期望，就是观测值 $y'$，加上因 $X$ 的变化 $x-x'$ 而预期在 $Y$ 上出现的变化。
\end{example}

\begin{example}
$e:\ X=x'$（处理对受处理者的效应，\S8.2.5）。
\begin{equation}
\begin{aligned}
E(Y_x\mid X=x')&=E(Y\mid x')+T(x-x')\\
&=rx'+T(x-x')\\
&=rx'+E\bigl(Y\mid\dox{x}\bigr)-E\bigl(Y\mid\dox{x'}\bigr),
\end{aligned} \tag{11.36}
\end{equation}
其中 $r$ 是 $Y$ 对 $X$ 的回归系数。
\end{example}

\begin{example}
$e:\ Y=y'$（例如：那些当前收入为 $Y=y'$ 的人，若我们强制要求每人每月接受 $x$ 小时培训，其预期收入 $Y$ 是多少）。
\begin{equation}
\begin{aligned}
E(Y_x\mid Y=y')&=y'+T\bigl[x-E(X\mid y')\bigr]\\
&=r'y'+E\bigl(Y\mid\dox{x}\bigr)-E\Bigl(Y\mid\dox\bigl(X=r'y'\bigr)\Bigr),
\end{aligned} \tag{11.37}
\end{equation}
其中 $r'$ 是 $X$ 对 $Y$ 的回归系数。
\end{example}

\begin{example}
考虑第 215 页的非递归价格—需求模型，即式 \textup{(7.9)}—\textup{(7.10)}：
\begin{equation}
q=b_1p+d_1i+u_1,\qquad p=b_2q+d_2w+u_2. \tag{11.38}
\end{equation}
我们的反事实问题（第 216 页）是：已知当前价格为 $P=p_0$，若我们把价格控制在 $P=p_1$，需求 $Q$ 的期望值将是多少？

作对应 $P\leftrightarrow X$、$Q\leftrightarrow Y$、$e\leftrightarrow\{P=p_0,i,w\}$，便可看出此问题与上面的例 11.7.2（处理对受处理者的效应）完全相同，只是要以 $i$ 与 $w$ 为条件。于是，由于 $T=b_1$，我们可以立即写出
\begin{equation}
\begin{aligned}
E(Q_{p_1}\mid p_0,i,w)&=E(Y\mid p_0,i,w)+b_1(p_1-p_0)\\
&=r_pp_0+r_ii+r_ww+b_1(p_1-p_0),
\end{aligned} \tag{11.39}
\end{equation}
其中 $r_p$、$r_i$、$r_w$ 分别是 $Q$ 对 $P$、$i$、$w$ 的回归中 $P$、$i$、$w$ 的系数。
\end{example}

式 \textup{(11.39)} 取代了第 217 页的式 \textup{(7.17)}。注意，价格方程 $p=b_2q+d_2w+u_2$ 的参数只经由回归系数进入 \textup{(11.39)}。因此，在它们由最小二乘直接估计的情形下，无须显式计算这些参数。

\begin{remarkenv}
例 11.7.1 其实并不令人惊讶；我们知道，在单调性假设下（第 293 页），因果概率是经验可识别的。但例 11.7.2 与例 11.7.3 引出如下猜想：
\end{remarkenv}

\noindent\textbf{猜想：}形如 $P(Y_x\mid e)$ 的任何反事实查询，在每个恒定效应模型中都是经验可识别的，即 $Y_{x_1}(u)-Y_{x_2}(u)$ 关于 $u$ 为常数。

以一个有挑战性的问题收尾是件好事。

\subsection{反事实的含义}

\qhead{问题（致作者）}

我很難理解反事实究竟有什么用。在我看来，它们似乎在回答一个错误的问题。在你的书里，你至少给出了几种不同的理由，说明何时需要反事实问题的答案，让我分别来谈：

\begin{enumerate}[label=\arabic*.,leftmargin=1.8em]
\item 责任的法律问题。从你的文字中我推断，美国法律体系认为：若被告造成了原告的损害，则被告有罪。但在我看来，法律显然有缺陷。责任应当系于被告行动的可预测后果，而不是实际发生了什么。一位医生为一种重病开出一种能治愈 99.99999\% 人群、却会致死 0.00001\% 人群的药物，即便他不幸而病人死亡，也不应被追究责任。如果法律建立在反事实的责任概念之上，那么在我看来，这法律有严重缺陷。
\item 决策中情境的使用。在第 217 页你写道：“此处值得强调，计算反事实期望的问题并非学术练习；事实上，它代表着几乎每一个决策情境中的典型情形。”我同意情境在决策中很重要，但不同意我们需要回答反事实问题。

在决策中，我们想估计的是 $P(\text{未来}\mid\dox{\text{行动}},\text{see}(\text{情境}))$。这当然是一个普通的 do 概率，而不是反事实查询。那么，我们为什么需要计算反事实？
\item 在书的后半部分，你用反事实来定义“$X$ 的原因”或“$Y$ 的充要原因”这类概念。同样，我理解人们忍不住要对这类日常语言中在用的概念作数学定义，但我不认为这通常很有帮助。我们为什么需要知道某一特定事件的“原因”？是的，我们感兴趣于知道事件的“原因”，因为它们使我们能预测未来，但这又回到上面第 2 点。
\end{enumerate}

把问题简化到最简形式，我的论证是：无论我们代表的是个人、企业、组织还是政府，我们都不断面对如何行动的决策（而这是我们仅有的决策！）。我们想知道的是：若我们以特定方式行动，最可能发生什么。所以我们想知道的是 $P(\text{未来}\mid\dox{\text{行动}},\text{see}(\text{情境}))$。我们并不想要、也不需要反事实的答案。

我的推理错在哪里？

\qhead{作者答复}

\begin{enumerate}[label=\arabic*.,leftmargin=1.8em]
\item 你的第一个问题质疑：在裁决法律责任时，使用单一事件概率而非总体概率是否明智。假定有一小部分病人对某种药物过敏，而制造商仍在不提示可能过敏反应的情况下销售该药。我们难道不认为：当一位过敏病人死亡时，他有权获得赔偿？通常，制药商会为这类例外情形投保，而不是在售药前让整个总体接受昂贵的检验——这在经济上划得来。医生当然被免责，因为他/她只是遵循了公认的医疗惯例。但法律确保：只要能证明（反事实地）若病人未服用该药，这一具体死亡本不会发生，就有人要为伤害付出代价。
\item 你的第二个问题涉及以观测结果为条件的决策。或者，用你的话说：“在决策中，我们想估计的是 $P(\text{未来}\mid\dox{\text{行动}},\text{see}(\text{情境}))$。”

问题在于，要正确地作出这一预测，我们需要一个按时间索引的序贯模型，其中“未来”变量与“情境”变量明确区分。然而，我们常常只被给予一个静态模型，其中“情境”变量被画成“行动”的后果；此时使用表达式 $P(y\mid\dox{x},z)$ 就不恰当了：它表示的是“我们先实施 $X=x$、随后观测到 $Z=z$”时 $Y=y$ 的概率，而所需要的是“我们先观测到 $Z=z$、然后实施 $X=x$”时 $Y=y$ 的概率。反事实给了我们表达所需概率的方式，即写作
\[
P\triangleq P(y_x\mid z),
\]
它代表：在“当前观测到 $Z=z$”的条件下，“若 $X$ 曾为 $x$，则 $Y=y$ 会发生”的概率。注意，若 $Z$ 在静态模型中是 $X$ 的非后代，则 $P\triangleq P(y_x\mid z)=P(y\mid\dox{x},z)$。

一个例子会有启发。设想一位工程师画出含一条门链 $X\to Y\to Z$ 的电路图 $M$。在时刻 $t_1$ 我们观测到 $Z(t_1)$，我们想知道以 $Z(t_1)$ 为条件的 $X(t_2)$ 对 $Y(t_3)$ 的因果效应。如果我们复制模型 $M$、为每个时隙 $t_i$ 指派一个模型 $M_i$（其中给出 $X(t_i)$、$Y(t_i)$、$Z(t_i)$ 之间、以及它们与此前变量的关系），然后计算 $P\bigl(y(t_3)\mid\dox{x(t_2)},z(t_1)\bigr)$，那么借助 do 演算、带上一切必要的时间指标，我们就能完成这一练习。但在静态模型里用反事实 $P(Y_x=y\mid z)$ 来做，要好得多。静态模型为我们省下了绘制等价序贯模型所需的千里长卷，而反事实使我们得以享受这一节省。

这是一项令人钦佩的发明。当然，有人可以主张：如果反事实断言不过是序贯模型等价物中科学预测的“会话简写”（第 118 页），那它们根本不需要。但这会违背科学方法的精神——在科学方法中，符号的经济性起着关键作用。原则上，代数中并不需要乘法——我们只靠加法也能活，把一个数自加所需次数即可。但若没有乘法，科学不会走得这么远。反事实亦然。
\item 归根到底，我们想了解“事件的原因”，确实是为了使我们能在未来更好地预测、更好地行动。但我们并不总是事先知道：在什么样的未来情境下，我们的知识会被召唤来帮助决策。我们需要知道某一具体事故之原因，理由可能是多方面的：提醒公众不要使用类似车辆、改进某些设备的维护程序、计算采取补救措施的成本、在敌方领土上造成类似事故，等等。每一种应用都可能从我们的模型中索取不同的信息，因此“事件的原因”可能是一种有用的方式，用以登记我们的经验，从而把它摊销到可能应用的广泛谱系上。
\end{enumerate}

\subsection{反事实的 d-分离}

\qhead{问题（致作者）}

我正试图把图 7.3 的孪生网络方法推广到反事实涉及两个以上可能世界的情形。考虑因果模型 $X\to Z\to Y$，并假定我们希望检验断言
\begin{equation}
Y_x\indep X\mid Y_z,Z_x,Y \tag{11.40}
\end{equation}
在该模型中是否成立。我会天真地构造如下的“三重网络”：

左部对应一个未施加任何干预的世界；中部对应施加 $\dox{X=x}$ 的世界；右部则施加 $\dox{Z=z}$。在这个网络中，\textup{(11.40)} 并不由 d-分离推出，因为从 $Y_x$ 到 $X$ 的路径因对 $Y$ 取条件而敞开。这种推广孪生网络方法的方式有什么问题吗？

\begin{figure}[htbp]
\centering
\begin{tikzpicture}[causal,scale=0.92,every node/.style={transform shape},font=\footnotesize]
  % 左世界
  \node[circle,fill,inner sep=1.1pt] (X) at (0,2.2) {};
  \node[font=\scriptsize,left=0.4mm] at (0,2.2) {$X$};
  \node[circle,fill,inner sep=1.1pt] (Z) at (0,1.2) {};
  \node[font=\scriptsize,left=0.4mm] at (0,1.2) {$Z$};
  \node[circle,fill,inner sep=1.1pt] (Y) at (0,0.2) {};
  \node[font=\scriptsize,left=0.4mm] at (0,0.2) {$Y$};
  \draw[ed] (X)--(Z); \draw[ed] (Z)--(Y);
  % 中世界 do(X=x)
  \node[font=\scriptsize] (xc) at (1.7,2.2) {$x$};
  \node[circle,fill,inner sep=1.1pt] (Zx) at (1.7,1.2) {};
  \node[font=\scriptsize,right=0.4mm] at (1.7,1.2) {$Z_x$};
  \node[circle,fill,inner sep=1.1pt] (Yx) at (1.7,0.2) {};
  \node[font=\scriptsize,right=0.4mm] at (1.7,0.2) {$Y_x$};
  \draw[ed] (xc)--(Zx); \draw[ed] (Zx)--(Yx);
  % 右世界 do(Z=z)
  \node[circle,fill,inner sep=1.1pt] (Xz) at (3.6,2.2) {};
  \node[font=\scriptsize,right=0.4mm] at (3.6,2.2) {$X_z$};
  \node[font=\scriptsize] (zc) at (3.6,1.2) {$z$};
  \node[circle,fill,inner sep=1.1pt] (Yz) at (3.6,0.2) {};
  \node[font=\scriptsize,right=0.4mm] at (3.6,0.2) {$Y_z$};
  \draw[ed] (zc)--(Yz);
  % 潜变量
  \node[circle,draw,inner sep=0.4pt,minimum size=2.2mm] (UX) at (1.8,3.3) {};
  \node[font=\scriptsize,above=0.2mm] at (1.8,3.35) {$U_X$};
  \node[circle,draw,inner sep=0.4pt,minimum size=2.2mm] (UZ) at (0.9,1.7) {};
  \node[font=\scriptsize,above=0.2mm] at (0.9,1.75) {$U_Z$};
  \node[circle,draw,inner sep=0.4pt,minimum size=2.2mm] (UY) at (1.7,-0.9) {};
  \node[font=\scriptsize,below=0.2mm] at (1.7,-0.95) {$U_Y$};
  \draw[edd] (UX)--(X); \draw[edd] (UX)--(Xz);
  \draw[edd] (UZ)--(Z); \draw[edd] (UZ)--(Zx);
  \draw[edd] (UY)--(Y); \draw[edd] (UY)--(Yx); \draw[edd] (UY)--(Yz);
\end{tikzpicture}
\caption{为检验式 \textup{(11.40)} 而画的三重网络。}
\label{fig:11.18}
\end{figure}

\qhead{作者答复（与 Ilya Shpitser 合作）}

你把孪生网络推广到两个以上世界，这是正确的；你的结论也正确：$Y_x$ 在给定 $Y_z,Z_x,Y$ 时与 $X$ 不独立。事实上，最近一篇论文（Shpitser and Pearl 2007）把反事实的这种多世界图形表示，系统表述于“反事实图”这一名目之下。

\qhead{追问（致作者）}

你的答复帮助我理解了孪生网络方法，也理解了《Direct and Indirect Effects》（Pearl 2001c）所用的方法；但它引出一个新问题：为什么在图 11.19 中，如 Pearl (2001) 所断言的那样，$Y_{xz}\indep Z_{x^*}\mid W$ 成立？若画一个“三重网络”，显然 $Y_{xz}$ 与 $Z_{x^*}$ 之间的路径并未被 $W$ 阻断；后者位于一个完全不同的世界。

\qhead{作者第二次答复}

独立性 $Y_{xz}\indep Z_{x^*}\mid W$ 在图 11.19 所示的图中确实成立。这是因为在你提到的“三重网络”中，$Y_{xz}$ 被变量 $W_{x^*}$ 与 $Z_{x^*}$ 分离。用 $W_{x^*}$ 替换 $W$ 的许可来自 do 演算的规则 3，它蕴含 $W_{x^*}=W$，因为 $X$ 是 $X$ 的非后代。这指出孪生网络推广中需要的一项重要改进：因果公理可能蕴含某些看似不同的反事实变量之间的等式约束，而在把 d-分离施用于反事实网络时，必须考虑这些隐藏的等式。Shpitser and Pearl (2007) 给出了编码与管理这些等式的一种系统方式。

\begin{figure}[htbp]
\centering
\begin{tikzpicture}[causal,scale=0.95,every node/.style={transform shape},font=\footnotesize]
  \node[circle,fill,inner sep=1.2pt] (X) at (0,1.6) {};
  \node[font=\small,left=0.5mm] at (0,1.6) {$X$};
  \node[circle,fill,inner sep=1.2pt] (Z) at (1.9,0.9) {};
  \node[font=\small,right=0.5mm] at (1.9,0.9) {$Z$};
  \node[circle,fill,inner sep=1.2pt] (W) at (3.7,1.5) {};
  \node[font=\small,right=0.5mm] at (3.7,1.5) {$W$};
  \node[circle,fill,inner sep=1.2pt] (Y) at (1.6,-0.6) {};
  \node[font=\small,below=0.5mm] at (1.6,-0.6) {$Y$};
  \draw[ed] (X)--(Z); \draw[ed] (Z)--(W); \draw[ed] (X)--(Y); \draw[ed] (Z)--(Y);
  % 潜变量
  \node[circle,draw,inner sep=0.4pt,minimum size=2.4mm] (U1) at (-1.3,0.5) {};
  \node[font=\scriptsize,left=0.3mm] at (-1.3,0.5) {$U_1$};
  \node[circle,draw,inner sep=0.4pt,minimum size=2.4mm] (U2) at (3.2,-0.3) {};
  \node[font=\scriptsize,right=0.3mm] at (3.2,-0.3) {$U_2$};
  \node[circle,draw,inner sep=0.4pt,minimum size=2.4mm] (U3) at (2.4,2.6) {};
  \node[font=\scriptsize,above=0.2mm] at (2.4,2.65) {$U_3$};
  \node[circle,draw,inner sep=0.4pt,minimum size=2.4mm] (U4) at (1.0,2.1) {};
  \node[font=\scriptsize,above=0.2mm] at (1.0,2.15) {$U_4$};
  \draw[edd] (U1) to[bend left=18] (X);
  \draw[edd] (U1) to[bend right=22] (Y);
  \draw[edd] (U2) to[bend left=20] (W);
  \draw[edd] (U2) to[bend right=18] (Y);
  \draw[edd] (U3) to[bend right=12] (X);
  \draw[edd] (U3) to[bend left=10] (W);
  \draw[edd] (U4) to[bend right=12] (X);
  \draw[edd] (U4) to[bend left=12] (Z);
\end{tikzpicture}
\caption{在其中，$Y_{xz}\indep Z_{x^*}\mid W$ 凭借等式 $W_{x^*}=W$ 而成立的图。}
\label{fig:11.19}
\end{figure}

\section{工具变量与不依从}

\subsection{不依从条件下的紧界}

\qhead{问题（致作者）}

我指的是你改进 Manski 在处理不完全依从时对处理效应之界的方式。你的方法利用了哪些 Manski 的方法没有利用的信息？其背后的直觉是什么？

\qhead{作者答复}

我们使用了与 Manski 相同的信息与相同的假设，并在由 $U$ 的典型划分所定义的 16 维空间中，用线性规划分析导出了紧界(Balke and Pearl 1994a, 1995a)。Manski 显然要么并未以求取紧界为目标，要么并未意识到把 $U$ 划分为其等价类的威力。回想一下，在 Frangakis and Rubin (2002) 以“主分层”(principal stratification)之名使之流行之前，这一划分对经济学家来说是未知的。

如我在第 269 页所述，Manski 的界在某些条件下是紧的，例如不存在“违拗者”(contrarian)之时。这意味着，只有当人群中存在违拗者时，人们才可能得到更窄的界，正如 Pearl (1995b) 所讨论的那些例子。那里表明，表明违拗者存在的数据如何能提供足够信息，使界塌缩为一点。该文也给出了这一现象如何发生的直观解释。

此处必须一提的是：典型划分的观念，配合 Balke and Pearl (1994a, 1995a,b) 所发展的线性规划方法，已在多种应用中成为强有力的分析工具。Tian and Pearl (2000) 把它用于界定因果概率；Kaufman 等 (2005) 与 Cai 等 (2008) 用它界定混杂中介存在时的直接效应；同样，Imai 等 (2008) 用它界定自然直接效应与自然间接效应。由该方法导出的闭式表达式，使研究者得以评估分布的哪些特征对于收窄界的宽度至关重要。

Rubin (2004) 在一项独立探索中，试图把典型划分应用于传统潜在结局框架内的直接效应与间接效应分析，但因缺乏图形与结构视角，他得出了这样的结论：这类效应是“定义不清的”“欺人更甚于助人”。我相信，本书的读者在潜在结局分析之结构根基的指引下，会得出更为积极的结论（见 \S4.5 与 \S11.4.2）。

\section{再论因果概率}

回首望去，在《因果性》首次出版八年之后，我认为第九章所得到的结果是反事实分析的一次胜利：它驱散了某些研究者就反事实之经验内容所表达的一切恐惧与惶惑(Dawid 2000; Pearl 2000)。它表明：一个乍看之下显得假设性的、定义不清的、不可检验的、因而配不上科学分析的量 $PN$，毕竟是可以定义、可以检验、并在某些情形下甚至可识别的。更何况，在某些数据组合之下、且不作任何假设，一个重要的法律主张——“若原告未服用该药，他本会活着”——竟能以概率一被确认，这实在令人惊叹。

\subsection{“以概率一有罪”究竟可能吗？}

我在会议与大学的几十场报告中展示过 \S9.3.4 的例子，而每次的结果都一样：式 \textup{(9.53)} 的 $PN\ge 1$ 总是招来普遍的不信——我们怎能从频率数据判定被告以概率一有罪，即 A 先生若未服用该药今天必定还活着？

Stephen Fienberg 教授听过其中两场报告，两次都摇着头说：“这不可能是真的。”在我看来，这一反应恰好构成一个坚实的证明：反事实分析能够产生非平凡的结果，因而它是真实的、有意义的、有用的；形而上学分析不会招致这样的抵触。

使人们不相信这一结果的，是该问题的三个令人困惑的方面：

\begin{enumerate}[label=\arabic*.,leftmargin=1.8em]
\item 一个假设性的、通常不可检验的量，竟能在某些条件下以概率一被确认；
\item 那些单独看来并未显示该药有实质效应的频率表，合在一起却蕴含对该药的完全易感性；
\item 一个未经检验的个体的属性，竟能依据从抽样总体取得的数据被赋予概率一。
\end{enumerate}

对认真对待逻辑与数学之益处的科学学生来说，第一个困惑其实并不令人惊讶。一旦我们赋予一个量以形式语义，我们就本质上定义了它与数据的关系；而某些条件下获得的数据足以把该量约束到可以精确确定的地步，并非不可设想。

反事实分析这一形式化带来的益处，其实可以在一个简单得多的例子里演示。考虑处理对受处理者的效应 $P(y_{x'}\mid x)$(\S8.2.5)。用文字说，$P(y_{x'}\mid x)$ 表示：一位受过处理($X=x$)的病人，若未曾受处理($X=x'$)，存活时间达到 $Y=y$ 的机会。这一反事实量似乎也抗拒经验测量，因为我们永远无法重演历史、拒绝对那些已受处理者施以处理。然而，对二值 $X$，我们可以写出
\[
P(y_{x'})=P(y_{x'}\mid x')P(x')+P(y_{x'}\mid x)P(x),
\]
并由此导出
\[
P(y_{x'}\mid x)=\bigl[P(y_{x'})-P(y,x')\bigr]/P(x).
\]
换言之，$P(y_{x'}\mid x)$ 可归约为经验可估计的量；$P(y_{x'})=P\bigl(y\mid\dox{x'}\bigr)$ 可在实验研究中估计，其余的量可在观测研究中估计。更何况，若数据支持等式 $P(y_{x'})=P(y,x')$，我们就可以稳妥地断定：一位受过处理的病人若未受处理，存活机会为零。那些先验地不信任反事实分析、视之为处理未定义量之演算的人，永远无缘享受这一发现：其中某些量竟是经验可定义的。逻辑在慷慨的时候，能为我们重演历史。

第二个困惑，已在式 \textup{(9.54)} 之后的段落中给出了直观解释。

第三个困惑才是让多数人震惊而不信的那个。尤其对统计学家而言，能对一个未经直接检验的具体个体说出任何确定的话，实属罕见。这源于两个因素。第一，统计学家通常处理有限样本，其变异性排除了任何断言中的确定性——不仅关于个体的断言，也包括关于底层分布之任何性质的断言。然而，这个因素不应进入我们的讨论，因为我们自始至终假定样本无限。（读者应设想表 9.2 中的数字代表百万量级。）

第二个因素源于这样一个事实：即便我们精确地知道一个分布，我们也不能对该分布中抽出的某个具体个体的属性赋予确定的概率估计。理由（据称）是：我们从不了解、更谈不上测量决定个体行为的一切解剖学与心理学变量；即便了解，我们也无法用手头分布所提供的粗糙类别把它们表示出来。于是，由于这种内在的粗糙性，“A 先生本会死亡”这句话永远不能被赋予概率一（事实上，不能被赋予任何确定的概率）。

这一由 Freedman and Stark (1999) 提出的论证，与日常话语中概率陈述的用法不相容，因为它蕴含：关于个体的每一句概率陈述，都必须是关于某个共享该个体全部特征的受限于总体的陈述。推到极端，这种限制性解读将坚持对原告作细密入微的刻画，而当一切相关细节都被计入时，就会把 $PN$ 化为零或一。很难设想这种解读会是司法标准之本意。立法者使用“更可能而非不然”的措辞，正是指示我们忽略那些没有数据可用的具体特征，而把判断建立在有可靠数据可用的最具体特征之上。在我们的例子里，A 先生的两项属性被注意到：(1) 他死了；(2) 他选择服用该药；这些在界定 $PN$ 时都被恰当地考虑在内。如果 A 先生的更多属性为人所知、并被认为相关（例如他红头发，或是左撇子），原则上也可以通过把分析限制在代表相应子总体的数据上来加以考虑。然而，在缺乏这类数据、且事先就知道我们永远无法匹配 A 先生的一切特异属性时，立法者的规定必须相对于手头的属性来解读。

\subsection{收窄因果概率的界}

Jin Tian 的系统工作(Tian and Pearl 2000)在若干方面改进了第九章的结果。第一，Tian 表明：对这些结果中的大多数而言，强外生性假设（式 \textup{(9.10)}）可以替换为弱外生性：
\[
Y_x\indep X\qquad\text{与}\qquad Y_{x'}\indep X.
\]
第二，在单调性假设（定义 9.2.13）下得到的估计量，在不假定单调性时构成下界。最后，第九章所导出的界是锐利的(sharp)，也就是说，若不加强假设便无法改进。

特别令人感兴趣的，是同时拥有实验研究与非实验研究数据、且不作其他任何假设时所得到的界。它们读作：
\begin{equation}
\begin{gathered}
\max\bigl\{0,\ P(y_x)-P(y_{x'}),\ P(y)-P(y_{x'}),\ P(y_x)-P(y)\bigr\}\ \le\ PNS\\
\le\ \min\bigl\{P(y_x),\ P(y'_{x'}),\ P(x,y)+P(x',y'),\ P(y_x)-P(y_{x'})+P(x,y')+P(x',y)\bigr\}
\end{gathered} \tag{11.41}
\end{equation}
\begin{equation}
\max\left\{0,\ \frac{P(y)-P(y_{x'})}{P(x,y)}\right\}\le PN\le \min\left\{1,\ \frac{P(y'_{x'})-P(x',y')}{P(x,y)}\right\} \tag{11.42}
\end{equation}
\begin{equation}
\max\left\{0,\ \frac{P(y_x)-P(y)}{P(x',y')}\right\}\le PS\le \min\left\{1,\ \frac{P(y_x)-P(x,y)}{P(x',y')}\right\} \tag{11.43}
\end{equation}

值得注意的是，在与药物有关的诉讼中，同时取得实验研究与观测研究的数据并不罕见。前者通常可在制造商或批准该药上市的机构（如 FDA）处获得，而后者则容易通过对总体的随机调查取得。在这类情形中，流行病学家用以确立法律责任的标准下界——超额风险比（式 \textup{(9.22)}）——可以通过使用式 \textup{(11.42)} 的下界而得到实质性改进。同样，式 \textup{(11.42)} 的上界可用于免除制药商的法律责任。Cai and Kuroki (2006) 分析了 $PNS$、$PN$ 与 $PS$ 的统计性质。

同样值得注意的是：在溯因任务中，即为最佳解释而推理时，$PN$ 是在事件 $Y=y$ 的诸竞争解释之间作抉择时最合理的度量。在此类应用中，\textup{(11.42)} 所给的界可由因果贝叶斯网络模型算出，其中 $P(y_x)$ 与 $P(y_{x'})$ 可经式 \textup{(9.48)} 计算。

\section*{致谢}
\addcontentsline{toc}{section}{致谢}

我感谢所有来信提问、并以此方式为本章作出贡献的读者。他们包括：David Bessler、Nimrod Megiddo、David Kenny、Keith A. Markus、Jos Lehmann、Dennis Lindley、Jacques A. Hagenaars、Jonathan Wilson、Stan Mulaik、Bill Shipley、Nozer D. Singpurwalla、Les Hayduk、Erich Battistin、Sampsa Hautaniemi、Melanie Wall、Susan Scott、Patrik Hoyer、Joseph Halpern、Phil Dawid、Sander Greenland、Arvid Sjolander、Eliezer S. Yudkowsky、UCLA 的 CS262Z（因果性研讨课，2006 年春季）学生，以及 UCLA 流行病学课程 EPIDEM 200C 的学员。

我同样感谢本书第一版的所有审稿人，以及帮助把这些书评提请读者注意的编辑们。他们包括：\textit{Choice}（Byerly 2000）、\textit{Structural Equation Modeling}（Shipley 2000a）、\textit{Chance}（McDonald 2001）、\textit{Technometrics}（Zelterman 2001）、\textit{Mathematical Reviews}（Lawry 2001）、\textit{Politische Vierteljahrsschrift}（Didelez and Pigeot 2001）、\textit{Technological Forecasting \& Social Change}（Payson 2001）、\textit{British Journal for the Philosophy of Science}（Gillies 2001）、\textit{Human Biology}（Chakraborty 2001）、\textit{The Philosophical Review}（Hitchcock 2001）、\textit{Intelligence}（O'Rourke 2001）、\textit{Journal of Marketing Research}（Rigdon 2002）、\textit{Tijdschrift Voor}（Decock 2002）、\textit{Psychometrika}（McDonald 2002b）、\textit{International Statistical Review}（Lindley 2002）、\textit{Journal of Economic Methodology}（Leroy 2002）、\textit{Statistics in Medicine}（Didelez 2002）、\textit{Journal of Economic Literature}（Swanson 2002）、\textit{Journal of Mathematical Psychology}（Butler 2002）、\textit{IIE Transactions}（Gursoy 2002）、\textit{Royal Economic Society}（Hoover 2003）、\textit{Econometric Theory}（Neuberg 2003）、\textit{Economica}（Abbring 2003）、\textit{Economics and Philosophy}（Woodward 2003）、\textit{Sociological Methods and Research}（Morgan 2004）、\textit{Review of Social Economy}（Boumans 2004）、\textit{Journal of the American Statistical Association}（Hadlock 2005），以及 \textit{Artificial Intelligence}（Kyburg 2005）。

我也感谢 UCLA 因果性博客（\url{http://www.mii.ucla.edu/causality/}）的各位撰稿人，以及该博客的维护者 William Hsu。

特别要向 Dennis Lindley 致以谢意与敬意：他不辞辛劳地从基本原理出发研读我的想法，使我得以断定，来自以统计学为基础的诸学科的读者将从本书中获益。我很幸运能与他的道路相交，并得以见证一位真正绅士的才智、好奇与正直。

若无 Jin Tian、Avin Chen、Carlo Brito、Blai Bonet、Mark Hopkins、Ilya Shpitser、Azaria Paz、Manabu Kuroki、Zhihong Cai、Kaoru Mulvihill 以及 UCLA 认知系统实验室全体成员的支持与洞见，本章不可能写成。过去六年里，当我被召唤去修补一个夺走我儿子 Daniel 生命的世界时（他于 2002 年在巴基斯坦卡拉奇被极端分子杀害），他们继续在因果性的青葱牧场上探索。这些年使我确信无疑：理性与启蒙的力量终将战胜狂热与不人道。

最后，我把这一版献给我的妻子 Ruth，感谢她在我们共同的磨难中给予的力量、爱与指引；献给我的女儿 Tamara 与 Michelle；献给我的孙辈 Leora、Torri、Adam、Ari 与 Evan，感谢他们在这段旅程中陪伴我，使之有意义、有目的。
